\section{The remaining terms}\label{sec:remaining}

The \emph{remaining terms} $\Gamma^*$ are the rational $\gamma$ with $\mathfrak g(\gamma)\ne0$, together with those that agree, have $\mathfrak g(\gamma)=0$ and move $A$. This section bounds their contribution
to the second moment (\cref{thm:remaining}). A block of $\Gamma^*$ is bounded, orbit by orbit under the unipotent elements $U^0_c$ that fix the frame, by Schur's test with a weight made of
volume ratios at the real places (\cref{thm:schur-orbit}); its \emph{size} depends on the cell rank and on a \emph{gauge} of the term. Squaring and summing over pairs of near terms, a pair gains
from the discriminant of its \emph{common lattice} in $X_0$ (\cref{thm:gain}); and lattices of given corank and small discriminant, and frames over them, are few
(\cref{thm:frame-count}). Pairs whose links differ are negligible; pairs with equal links gain the norm of the places where the block meets.

\emph{Conventions.} Throughout the data are admissible (\cref{ss:filling-data}) and $E$ contains a primitive $5040$th root of unity. We write $f$ for the number of infinite places of $E$, $k=\dim X_0$,
$n=\varrho+1$, $\tau=b^{kf/4}$, so that $\tau^4=b^{kf}$, and $n_p=\#\bar R_0\,n_XN^{-kf}$, with $n_X$ the number of classes. Since $k\ge1000(2\varrho+2)^9$, every power $b^e$ with $e\le10^5n^9f$
is at most $\tau$, and so is the number $n^3+1$. For a finite idele $a$ the \emph{generalised index} $[a]$ is the inverse of the factor by which multiplication by $a$ scales Haar measure on the finite
adeles, and $[a]=0$ if $a$ is not a unit; for integral $a$ it is $\#\hat\OO/a\hat\OO$, and for $a\in E^\times$ it is $\mathrm{Nm}(a)=\abs{\mathrm N_{E/\Q}(a)}=\prod_v\abs a_v^2$, the product over the infinite places with the
chosen embeddings. For the discriminant $\mathfrak g(\gamma)$, which is fixed by complex conjugation, $\mathfrak g_{\rm Nm}=\prod_v\abs{\mathfrak g(\gamma)}_v$; it is a positive integer when $\mathfrak g(\gamma)\ne0$.
The lattice of a class $c$ is written $\Lambda$, and vectors are written in its coordinates, in which $h_c^{-1}\gamma h_{c'}$ is the matrix of a term.
\lean{New.genIndex, New.submoduleIndex, New.colLattice, New.Filling.finAdeleMatrix, New.Filling.finAdeleMatrix_conjTranspose, New.Filling.finAdeleMatrix_smul, New.Filling.map_conjTranspose, New.Filling.finAdeleMatrix_mul, New.Filling.finAdeleMatrix_add}

\subsection{Lattices over the finite adeles}\label{ss:lattices}
Let $\hat\OO$ be the ring of integral finite adeles of $E$. A \emph{lattice of rank $r$} in $\Ad_f^\kappa$ is the $\hat\OO$-span $L$ of the columns of $xa$, where $x$ is a $\kappa\times r$ matrix over $E$ with
independent columns and $a\in\GL_r(\Ad_f)$; its \emph{discriminant} is $\disc L=[\det((xa)^*(xa))]^{1/2}$. For a matrix $H$ over $\Ad_f$ with a left inverse (the passage to the coordinates of a lattice) and
a subspace $V$ of $E^\kappa$, the \emph{saturated lattice} $\Lambda_V$ is the set of points of the adelic span of $V$ that $H$ makes integral; its discriminant is written $d_H(V)$.
\lean{New.Filling.IsLatticeOfRank, New.Filling.discriminant, New.Filling.saturatedLattice, New.Filling.sublatticeBasis, New.Filling.sublatticeDiscriminant, New.Filling.sublatticeIndex}

\begin{lemma}[Discriminants]\label{lem:lattice-disc}
\begin{enumerate}[label=\textup{(\alph*)}]
\item The discriminant does not depend on the basis $xa$ of $L$, and it is positive.
\item If $M\subseteq L$ is an $\hat\OO$-submodule of finite index, then $M$ is a lattice of rank $r$ and $\disc M=[L:M]\disc L$.
\item A matrix $\eta$ over $E$ with a left inverse maps lattices of rank $r$ to lattices of rank $r$; if $\eta^*\eta=1$, it keeps the discriminant.
\item $\Lambda_V$ is a lattice of rank $\dim V$, and $d_H(V)\ge1$ if every adelic matrix $m$ with $Hm$ integral has an integral Gram matrix $m^*m$. If $H$ maps a lattice $L$ of rank $r$ into
integral vectors, then for the span $W$ of $L$ one has $L\subseteq\Lambda_W$ with finite index $i$, and $\disc L=i\,d_H(W)$.
\item Every ideal of finite index of $\hat\OO$ is principal.
\end{enumerate}
\end{lemma}
\lean{New.Filling.genIndex_det_gram_eq_of_colLattice_eq, New.Filling.discriminant_colLattice, New.Filling.IsLatticeOfRank.discriminant_pos, New.Filling.det_gram_ne_zero, New.Filling.hasLeftInv_finAdeleMatrix_mul, New.Filling.exists_colLattice_fromCols_eq, New.Filling.exists_colLattice_eq_of_submoduleIndex_ne_zero, New.Filling.sqrt_genIndex_det_gram_mul, New.Filling.exists_colLattice_mul_eq_of_le, New.Filling.IsLatticeOfRank.of_submoduleIndex_ne_zero, New.Filling.IsLatticeOfRank.discriminant_eq_mul_submoduleIndex, New.Filling.IsLatticeOfRank.of_le_and_discriminant_eq_mul_index, New.Filling.colLattice_le_preimageLattice_iff, New.Filling.map_colLattice_finAdeleMatrix_mul, New.Filling.IsLatticeOfRank.map_of_hasLeftInv, New.Filling.IsLatticeOfRank.discriminant_map_of_unitary, New.Filling.IsLatticeOfRank.map_of_unitary, New.Filling.saturatedLattice_le, New.Filling.isLatticeOfRank_saturatedLattice, New.Filling.discriminant_saturatedLattice, New.Filling.one_le_sublatticeDiscriminant, New.Filling.IsLatticeOfRank.exists_eq_colLattice_sublatticeBasis_mul, New.Filling.isPrincipal_of_cardQuot_ne_zero, New.Filling.colLattice_mul, New.Filling.exists_isIntMatrix_mul_eq, New.Filling.colLattice_mul_le, New.Filling.exists_range_eq_colLattice, New.IsValuationPair.exists_square_colLattice_eq, New.IsValuationPair.isBezout, New.Filling.exists_isPrimitiveMatrix_mul_finAdeleMatrix, New.Filling.isPrimitiveMatrix_mul_sublatticeBasis, New.Filling.exists_sublatticeBasis_eq_finAdeleMatrix_mul, New.Filling.exists_finAdeleMatrix_mul_eq_sublatticeBasis_mul, New.Filling.sublatticeDiscriminant_sq_mul_genIndex_sq, New.Filling.sublatticeDiscriminant_pos, New.Filling.sublatticeDiscriminant_mul_index_sq_le}

\begin{proof}
(a) Two adelic bases $B$, $C$ of the same $\hat\OO$-module, $B$ with a left inverse, satisfy $C=BU$ and $B=CU'$ with $U$, $U'$ integral, so $U$ lies in $\GL_r(\hat\OO)$ and $[\det U]=1$. Since
$\det(C^*C)=\overline{\det U}\det U\det(B^*B)$ and the generalised index is multiplicative, the two Gram determinants have the same index. Positivity: at a complex embedding the Gram matrix of $r$
independent vectors is positive definite, so $\det(x^*x)\ne0$, and $[\det((xa)^*(xa))]=\mathrm{Nm}(\det x^*x)[\det a]^2>0$.

(b) Over each completion $\hat\OO_v$ is a principal ideal domain, and $\hat\OO$ satisfies the following exchange property: any $a$, $b$ are multiples of $(1-e)a+eb$ for an idempotent $e$. Column operations over $\hat\OO$
therefore bring a generating matrix $(C\mid R)$ of a module of finite index in $\hat\OO^r$ to a square matrix; hence $M$ has a basis $BA$ with $A$ integral and invertible over $\Ad_f$, and
$[L:M]=[\hat\OO^r:A\hat\OO^r]=[\det A]$. Then $\disc M^2=[\det A]^2\disc L^2$.

(c) The image of the span of the columns of $xa$ is the span of the columns of $(\eta x)a$, and $\eta x$ has independent columns. If $\eta^*\eta=1$ the Gram matrix does not change.

(d) There are a basis $y$ of $V$ and $a\in\GL_s(\Ad_f)$, $s=\dim V$, such that $Hya$ is primitive, that is integral with an integral left inverse $p$ (take a basis of the integral points of the span locally at each
place); the columns of $m=ya$ span $\Lambda_V$, because a point $mt$ of $\Lambda_V$ has $t=p(Hmt)$ integral. If $Hm$ integral forces $m^*m$ integral, the Gram determinant of $m$ is a nonzero integral idele, whose
index is at least $1$. Finally let $L$ have the basis $xa$ and let $W$ be the span of $x$, with chosen basis $S=yb$ of $\Lambda_W$. Then $x=yT$ with $T$ invertible over $E$, and $xa=SA$ with $A=b^{-1}Ta$; the
matrix $A=p\,H(xa)$ is integral, so $L\subseteq\Lambda_W$ with index $[\det A]$, and $\disc L=[\det A]\,d_H(W)$ as in (b).

(e) An ideal $I$ of finite index contains its index, a nonzero integer $m$. Hence $I$ is determined by its images in the finitely many factors $\hat\OO_v$ with $v\mid m$, each a power of the maximal ideal,
and the idele that is a uniformiser to that power at these places and $1$ elsewhere generates it.
\end{proof}

For an $R$-module $X$ let $\zeta_X(s)=\sum_M[X:M]^{-s}$, the sum over the submodules of finite index.
\lean{New.Filling.submoduleZeta, New.Filling.submoduleZeta_anti, New.Filling.cardQuot_comap, New.Filling.cardQuot_comap_of_surjective, New.Filling.tsum_relIndex_range_le_submoduleZeta, New.Filling.submoduleZeta_le_of_surjective, New.Filling.cardQuot_eq_cardQuot_map_fst_mul_comap_inr, New.Filling.eq_span_sup_of_map_fst_eq_span, New.Filling.tsum_quotient_cardQuot_inv_pow}

\begin{lemma}[Sums over sublattices]\label{lem:sublattice-zeta}
Let $R$ be a commutative ring whose ideals of finite index are principal.
\begin{enumerate}[label=\textup{(\alph*)}]
\item $\zeta_{R\times N}(s+1)\le\zeta_R(s+1)\,\zeta_N(s)$ for every $R$-module $N$.
\item $\zeta_{R^r}(r+\mu)\le\zeta_R(\mu)^r$.
\item For a lattice $L$ of rank $r$ over $\hat\OO$, $\sum_M[L:M]^{-(r+\mu)}\le\zeta_{\hat\OO}(\mu)^r$. For $\mu=2\varrho+4$ one has $\zeta_{\hat\OO}(\mu)\le2$, so the sum is at most $2^r\le2^k$ when $r\le k$.
\end{enumerate}
\end{lemma}
\lean{New.Filling.submoduleZeta_prod_le, New.Filling.submoduleZeta_pi_le, New.Filling.tsum_submoduleIndex_colLattice_inv_pow_le, New.Filling.IsLatticeOfRank.tsum_submoduleIndex_inv_pow_le, New.Filling.tsum_index_saturatedLattice_le}

\begin{proof}
(a) Let $M\subseteq R\times N$ have finite index. Its first coordinates form an ideal $I=(g)$ of finite index, and with $M'=M\cap N$ and any $n$ with $(g,n)\in M$, the module $M$ is spanned by $(g,n)$
and $M'$; the class of $n$ modulo $M'$ determines $M$. The index is $[R\times N:M]=[R:I]\,[N:M']$, and there are $[N:M']$ classes. Hence
\[
  \zeta_{R\times N}(s+1)\le\sum_I\sum_{M'}[N:M']\,\bigl([R:I][N:M']\bigr)^{-(s+1)}=\zeta_R(s+1)\,\zeta_N(s).
\]
(b) Induction on $r$, with $R^r=R\times R^{r-1}$: by (a), $\zeta_{R^r}(r+\mu)\le\zeta_R(r+\mu)\zeta_{R^{r-1}}(r-1+\mu)\le\zeta_R(\mu)\cdot\zeta_R(\mu)^{r-1}$, since $\zeta_R$ decreases in $s$.
(c) A lattice of rank $r$ is isomorphic to $\hat\OO^r$ through its basis, and the ideals of finite index of $\hat\OO$ are principal (\cref{lem:lattice-disc}(e)). Contraction to $\OO_E$ is injective on ideals
of finite index and keeps the index, so $\zeta_{\hat\OO}(\mu)$ is at most the Dedekind zeta value $\zeta_E(\mu)$, an Euler product; for $\mu=2\varrho+4$ and $(\varrho+1)[E:\Q]\le2^{2\varrho}$ its logarithm
is at most $(\log2)/(\varrho+1)$, so $\zeta_{\hat\OO}(\mu)^{\varrho+1}\le2$ and $\zeta_{\hat\OO}(\mu)\le2$.
\end{proof}

Fix $u$ in the finite adelic unitary group of $X_0=E^\kappa$ and consider the unimodular lattice $u\hat\OO^\kappa$ and the rational lattice $E^\kappa\cap u\hat\OO^\kappa$. For a subspace $T$ let $d(T)$ be the
discriminant of $T\cap u\hat\OO^\kappa$, a saturated lattice for $H=u^{-1}$, and write $\norm X$ for the Euclidean length of a vector over all infinite places together; $\mathrm{rd}_E$ is the root
discriminant.

\begin{proposition}[Siegel's lemma]\label{prop:siegel-lemma}
\begin{enumerate}[label=\textup{(\alph*)}]
\item (Jacobi.) $d(T)=d(T^\perp)$ for every subspace $T$.
\item (Siegel's lemma.) A subspace $T$ of codimension $s$ is the orthogonal of $s$ nonzero vectors $X_1,\dots,X_s$ of $E^\kappa\cap u\hat\OO^\kappa$ with
\[
  \Bigl(\prod_i\norm{X_i}\Bigr)^{2f}\le(8fs\,\mathrm{rd}_E)^{fs}\,d(T).
\]
\item $\sum_Td(T)^{-(k+1)}\le(200fk\,\mathrm{rd}_E)^{f(k+1)s}$, the sum over the subspaces $T$ of codimension $s$.
\end{enumerate}
\end{proposition}
\lean{New.Filling.siegel_lemma, New.Filling.tsum_sublatticeDiscriminant_inv_le, New.Filling.genIndex_sq_mul_abs_norm_eq, New.Filling.exists_det_ne_zero_prod_frameLength_le, New.Filling.one_le_ideleAbs_det_mul_sqrt, New.Filling.norm_archCoords_mulVec, New.Filling.frameLength_mkAdelicGL, New.Filling.mulVec_mem_baseLattice, New.Filling.one_le_norm_archCoords, New.Filling.exists_inv_sublatticeDiscriminant_pow_le_prod, New.Filling.finrank_ker_conjTranspose, New.Filling.ker_conjTranspose_mul_transpose, New.Filling.jacobi_perp_discriminant, New.Filling.IsKernelExtension, New.Filling.exists_isKernelExtension, New.Filling.mul_eq_zero_of_colSpan_le_ker, New.Filling.IsKernelExtension.colSpan_fromCols, New.Filling.IsKernelExtension.colSpan_le, New.Filling.IsKernelExtension.eq_zero_of_mulVec_mem, New.Filling.IsKernelExtension.finrank_eq, New.Filling.IsKernelExtension.hasLeftInv, New.Filling.IsKernelExtension.mul_eq_zero}

\begin{proof}
(a) Let $xb$ and $ya$ be the chosen bases of $T\cap u\hat\OO^\kappa$ and $T^\perp\cap u\hat\OO^\kappa$; the matrices $u^{-1}xb$ and $u^{-1}ya$ are primitive and $x^*y=0$. In the unimodular lattice $\hat\OO^\kappa$ two
orthogonal primitive sublattices of complementary ranks have Gram determinants of the same index: complete $u^{-1}xb$ to a basis of $\hat\OO^\kappa$; the inverse of its Gram matrix is integral and its block on the
complement is the inverse of the Gram matrix of the dual basis, which spans the projection of $u^{-1}ya$; this is Jacobi's identity for complementary minors of a unimodular matrix. As $u^{-1}$ is unitary it drops
out of the Gram matrices, which gives $[\det a]^2\,\mathrm{Nm}(\det y^*y)=[\det((xb)^*(xb))]$.

(b) Take $y$ and $a$ as in (a) for $T^\perp$, which has dimension $s$, and at each infinite place $v$ an $r_v\in\GL_s(\C)$ with $(y_vr_v)^*(y_vr_v)=1$. The adelic matrix $(r,a)$ has
$\abs{\det(r,a)}\cdot d(T)\ge1$: its infinite part cancels the norm of $\det y^*y$, and its finite part the index of $\det a$, by (a). Minkowski's second theorem over $E$ \cite{BombieriVaaler1983}, in its adelic
form for $(r,a)$, gives an invertible $B$ over $E$ with $a^{-1}B^{\rm T}$ integral and
\[
  \abs{\det(r,a)}\prod_i\ell(B_i)^{2f}\le(2fs\,\mathrm{rd}_E)^{fs},
\]
where $\ell(B_i)$ is the Euclidean length of the vectors $r_v^{-1}B_i$ at all places together. Put $X_i=yB_i$. Then $u^{-1}yB^{\rm T}=(u^{-1}ya)(a^{-1}B^{\rm T})$ is integral, so $X_i$ lies in $E^\kappa\cap u\hat\OO^\kappa$; $X_i\ne0$
because $B$ is invertible and $y$ has a left inverse; the $X_i$ span $T^\perp$, whose orthogonal is $T$; and $\norm{X_i}=\ell(B_i)$, because $y_vr_v$ is an isometry. Multiplying by $\abs{\det(r,a)}d(T)\ge1$ gives the bound,
with $2\le8$.

(c) Raise (b) to the power $k+1$: with $m=f(k+1)$ and $C=(8fk\,\mathrm{rd}_E)^m$, $d(T)^{-(k+1)}\le\prod_i C\norm{X_i}^{-2m}$. The subspace $T$ is recovered from the tuple $(X_i)$ as its orthogonal, so the sum over
$T$ is at most the $s$th power of $C\sum_X\norm X^{-2m}$ over the nonzero lattice vectors. These have length at least $1$ (an algebraic integer whose conjugates are all smaller than $1$ in modulus is $0$), so
they are $1$-separated, and since $2m$ exceeds the real dimension $2fk$ by at least $2$, a packing count over the shells $2^t\le\norm X<2^{t+1}$ gives $\sum_X\norm X^{-2m}\le5^{2m}$. As $200=5^2\cdot8$, this is the bound.
\end{proof}

\subsection{Schur's test for one orbit}\label{ss:schur-orbit}
For $\gamma\in G_1$ and $y$, let $\nu_\gamma(y)$ be the ratio of the volumes of the projections of $\gamma f$ and of $f$ to the negative space $yV_-$ of $J_1$. A term of cell rank $j$ has $i=\min(\varrho,j)$, and
its \emph{Schur weight} at $y$ is the boxed cutoff at $\gamma^{-1}y$ times $\prod_v\nu_\gamma(y_v)^{-(k+i)}$.
\lean{New.frameVolume, New.volumeRatio, New.volumeRatio_pos, New.volumeRatio_inv, New.volumeRatio_mul_left, New.measurable_volumeRatio, New.continuous_frameVolume, New.frameVolume_pos, New.frameVolume_mul_left, New.Filling.boxCutoff, New.Filling.besselCutNormSq, New.Filling.schurWeight, New.Filling.schurIntegral, New.Filling.measurable_boxCutoff, New.Filling.measurable_besselCutNormSq, New.Filling.measurable_schurWeight}

\begin{lemma}[Boxed points and integral vectors]\label{lem:box-integral}
\begin{enumerate}[label=\textup{(\alph*)}]
\item If $A$ maps injectively to vectors in $E^\iota$ whose entries are algebraic integers of modulus at most $r$ at every infinite place, then $\#A\le(1+2r)^{2\#\iota f}$.
\item The boxed points form a measurable set, invariant under left multiplication by unipotent elements fixing every frame vector at every place. At a boxed point $y$ the entries of the
Gram--Schmidt matrix have modulus at most $b^{26}$, and $y_v$ lies in the frame box of radius $b^{27}\ell(y_v)\le b^{37n}$.
\end{enumerate}
\end{lemma}
\lean{New.Filling.encard_le_of_isIntegral_norm_le, New.Filling.finite_setOf_levelMultiples_add_le, New.Filling.mem_levelMultiples, New.Filling.mem_boxedPoints_of_reducedIndicator_ne_zero, New.Filling.mul_mem_reducedPoints_iff, New.Filling.norm_basisScale_div_le_one, New.Filling.ratTriangle_eq_one_apply_of_le, New.exists_measurableSet_existsUnique_add_mem}

\begin{proof}
(a) Embed $A$ into $\C^{\iota\times\{v\}}$, a real space of dimension $2\#\iota f$. The image lies in the ball of radius $r$ for the maximum norm, and distinct points are at distance at least $1$, because a
nonzero algebraic integer has a conjugate of modulus at least $1$. Balls of radius $1/2$ around the points are disjoint and lie in the ball of radius $r+1/2$, which gives the count. (b) The reduced points are
defined by the entries above the diagonal of the triangular matrix, a measurable function, lying in a measurable set $D$ (\cref{prop:packing}(c)); such a $D$ exists because in each coset of $M\OO_E$ a measurable
injection to $\R$ picks a least point of a bounded set meeting the coset in a finite nonempty set. A unipotent $s$ fixing the frame leaves the frame vectors $y^{-1}f_a$ and the triangular matrix unchanged, and
it does not change the length. The scale of the frame vector of index $a$ is $\kappa_v^{\varrho-a}/\sqrt N$ with $\abs{\kappa_v}\ge1$, so ratios of later to earlier scales are at most $1$; a frame vector is a sum of
at most $\varrho+1\le b$ orthogonal vectors, each not longer than $\ell(y_v)$, with coefficients of modulus at most $b^{26}$. Hence its norm is at most $b^{27}\ell(y_v)$, and in the cut $b^{27}b^{6n+2}\le b^{37n}$.
\end{proof}

\begin{lemma}[A triangular basis of the negative space]\label{lem:triangular-basis}
Let $\gamma,y\in G_1$, let a norm-one pair occur for $\gamma$ from $j$, $i=\min(\varrho,j)$, and let the block $B$ of $Z(\gamma)$ on the subframe $e_1,\dots,e_i,z$ be invertible. There are vectors
$\omega_0,\dots,\omega_i$ of $yV_-$ with Gram matrix $-1$ whose products with the subframe form a lower triangular invertible matrix $L$ and whose products with the frame are those of the basis of the span
of $\gamma e_1,\dots,\gamma e_i,\gamma z$ dual to the subframe, times $L$. For such $(\omega,L)$:
\begin{enumerate}[label=\textup{(\alph*)}]
\item $\omega$ is a basis of $yV_-\cap I_\gamma^\perp$ pairing perfectly with the subframe, $I_\gamma=A\cap\gamma A$;
\item $\abs{\det L}\le\rho^{i+1}$ if the frame vectors of $y$ have norm at most $\rho$, and $\abs{\det L}=\abs{\det L_1}\abs{L_{zz}}$ for the leading block $L_1$;
\item with $\omega'=\gamma(e_1,\dots,e_i,z)B^{-1}L$ and $v$ the parts in $X_0$ of the columns of $\omega'-\omega$, the differences $\omega'-\omega$ are orthogonal to the frame and have the Gram matrix of $v$;
\item if $y$ has length at most $\theta$, then $\abs{L_{zz}}^2\le\theta^{2(j-i)}(1+\abs{v_z}^2)$, with equality $\abs{L_{zz}}^2=1+\abs{v_z}^2$ when $j\le\varrho$;
\item $\abs{\det B}=\abs{D_j}$ for the leading minor $D_j$ of order $j$ of $Z(\gamma)$.
\end{enumerate}
\end{lemma}
\lean{New.RealPlace.exists_isTriangularBasis, New.RealPlace.IsTriangularBasis.isNegativeBasis, New.RealPlace.IsTriangularBasis.norm_det_le, New.RealPlace.IsTriangularBasis.norm_det_eq_mul_norm_last, New.RealPlace.IsTriangularBasis.conjTranspose_frameCols_mul_J₁_mul_sub_eq_zero, New.RealPlace.IsTriangularBasis.gram_differenceX₀, New.RealPlace.IsTriangularBasis.norm_last_le, New.RealPlace.IsTriangularBasis.norm_last_sq_eq_one_add_dotProduct_differenceX₀, New.RealPlace.IsTriangularBasis.norm_last_sq_le, New.RealPlace.IsTriangularBasis.apply_eq_frameVector_dotProduct, New.RealPlace.IsTriangularBasis.norm_dotProduct_J₁_mul_col_le_norm, New.RealPlace.IsTriangularBasis.conjTranspose_subframeCols_mul_J₁_mul_eq, New.RealPlace.IsTriangularBasis.det_ne_zero, New.RealPlace.IsTriangularBasis.gram_eq_neg_one, New.RealPlace.IsTriangularBasis.products_eq, New.RealPlace.IsTriangularBasis.triangular, New.RealPlace.IsTriangularBasis.J₁_mul, New.RealPlace.IsNegativeBasis.J₁_mul, New.RealPlace.IsNegativeBasis.det_ne_zero, New.RealPlace.IsNegativeBasis.orthogonal, New.RealPlace.inclMatrix_mul_dualProducts, New.RealPlace.frameCols_mul_mul_subframeCols, New.RealPlace.subframeIndex, New.RealPlace.subframeIndex_castSucc, New.RealPlace.subframeIndex_last, New.RealPlace.subframeIndex_injective, New.RealPlace.subframeIndex_self, New.RealPlace.le_and_ne_top_of_notMem_range, New.RealPlace.J₁_mulVec_eq_self, New.RealPlace.form_frameCols_mulVec, New.RealPlace.posSemidef_neg_gram_mul_frameCols, New.RealPlace.gram_definiteCols_mul_eq, New.RealPlace.J₁_mul_coe, New.RealPlace.J₁_mul_one_sub_J₁_mul, New.RealPlace.J₁_mul_inv_conjTranspose, New.RealPlace.J₁_mul_frameLeftInverse_conjTranspose, New.RealPlace.exists_eq_one_sub_J₁_mul, New.RealPlace.eq_mul_of_J₁_mul_eq_neg, New.RealPlace.frameLeftInverse, New.RealPlace.frameLeftInverse_mul_frameCols, New.RealPlace.frameCols_J₁_gram, New.RealPlace.realFrameGram_eq, New.RealPlace.conjTranspose_realFrameGram_mulVec}

\begin{proof}
The space $V_-$ is spanned by the columns of $(1-J_1)f$, and $yV_-$ is negative definite of dimension $\varrho+1$; a vector of $yV_-$ is determined by its products with the frame $f$, since the frame pairs perfectly
with $yV_-$ (the form is $\le0$ on the span of $f$, and $-\abs{x_z}^2$ on $fx$). Let $\omega_0$ be the vectors of $yV_-$ whose products with the frame are those of the dual basis
$\gamma(e_1,\dots,e_i,z)B^{-1}$; Gram--Schmidt for the negative form applied to $\omega_0$ gives $\omega=\omega_0L$ with $L$ lower triangular, invertible, and Gram matrix $-1$; the products of the subframe with
$\omega$ form $L$, because those with $\omega_0$ form the identity.

(a) A vector of $I_\gamma$ lies in $\gamma A$, so it is orthogonal to the frame of $\gamma$, hence to $\omega$; and the products with the subframe form the invertible $L$. As $\dim yV_-\cap I_\gamma^\perp=\varrho+1-\dim I_\gamma=i+1$,
$\omega$ is a basis. (b) $L_{al}$ is the product of $y^{-1}$ applied to the $a$th subframe vector with the unit vector $y^{-1}\omega_l$ of $V_-$, so $\abs{L_{al}}\le\rho$; Hadamard's inequality on the triangular $L$ gives
$\abs{\det L}\le\prod_a\abs{L_{aa}}\le\rho^{i+1}$, and the determinant of a triangular matrix is the product of the leading block and the last entry.
(c) The products of $\omega'$ and of $\omega$ with the frame coincide by construction, so $\omega'-\omega$ is orthogonal to the frame; a vector orthogonal to the frame has vanishing coordinates at the $e^*_a$
and at $z$ and its form is that of its part in $X_0$, up to the isotropic coordinates, which pair to zero.
(d) If $j\le\varrho$ the flags of $f$ and $\gamma f$ share a vector $fx=\gamma f x'$ with $x'_z=1$ and $\abs{x_z}=1$; it is a combination of the image of the subframe, and $\omega'_z$ is a multiple of it, hence lies in the span of
$f$, where it is orthogonal to the difference $\omega'_z-\omega_z$ (which is orthogonal to the frame). Therefore
\[
  -1=\langle\omega_z,\omega_z\rangle=\langle\omega'_z,\omega'_z\rangle+\abs{v_z}^2=-\abs{L_{zz}}^2+\abs{v_z}^2.
\]
At the top level $i=\varrho$ take $u$ unipotent with $\norm{uy}\le\theta$; as $\omega_z$ is orthogonal to $A$, its product with $z$ is that with $u^{-1}z$, so $\abs{L_{zz}}\le\norm{uy}\le\theta$, and $1+\abs{v_z}^2\ge1$.
(e) If $j=\varrho+1$ the subframe is the whole frame and $B=Z(\gamma)$. Otherwise $i=j$, and the shared vector at the index of $z$, corrected by the shared vectors at the indices from $i$ to $\varrho-1$, is a vector
$fx=\gamma fx'$ with $x'$ supported on the subframe, $x'_z=1$ and $\abs{x_z}=1$; its products with the subframe are $(0,\dots,0,-x_z)$, because the $e_a$ are isotropic and orthogonal to every frame vector, and $z$ has norm $-1$
for $J_1$. By Cramer's rule, replacing the last column of $B$ by $Bx'$ leaves the determinant unchanged, since $x'_z=1$; expanding along this column, $\abs{\det B}$ is the modulus of the leading
minor of order $i$.
\end{proof}

\begin{proposition}[The weight at one real place]\label{prop:volume-ratio}
Let a norm-one pair occur for $\gamma_0$ from $j$, $i=\min(\varrho,j)$, $Z$ the frame products of $\gamma_0$ with $\delta=\det$ of its block on the subframe $e_1,\dots,e_i,z$ nonzero, and let $y$ have length at most $\theta$
and frame vectors of norm at most $\rho$, $k\ge7$.
\begin{enumerate}[label=\textup{(\alph*)}]
\item $\nu_\gamma(y)$ is the modulus of the determinant of the products of a negative basis of $yV_-\cap I_\gamma^\perp$ with $\gamma e_1,\dots,\gamma e_i,\gamma z$, divided by that with $e_1,\dots,e_i,z$; and
$\nu_\gamma(y)\ge2^{-i-1}\abs\delta\det(1+vv^*)/\abs{\det L}^2$ for a triangular basis $(\omega,L)$, with $v$ the parts in $X_0$ of the differences of two bases, which move affinely when $y$ is moved by a
unipotent element fixing the frame.
\item There is a measurable $\phi$ on $X_0^i$ with $\int\phi\le(2\rho^2)^{(i+1)i}(2^{i+1}\theta^{2(j-i)})^k/\abs\delta^{k+i}$ and $\nu_\gamma^{-(k+i)}\le\phi(A)$ whenever $\gamma$ has the same $Z$ and norm-one pairs from $j$, and
its subframe differs in $X_0$ from that of $\gamma_0$ by ${}^t\bar A$ times the first $i$ rows of the block.
\end{enumerate}
\end{proposition}
\lean{New.RealPlace.IsNegativeBasis.volumeRatio_eq, New.RealPlace.IsTriangularBasis.norm_det_mul_le_volumeRatio, New.RealPlace.differenceX₀_eq_add, New.RealPlace.exists_majorant_volumeRatio_inv_pow, New.RealPlace.IsNegativeBasis, New.RealPlace.IsTriangularBasis, New.RealPlace.exists_isTriangularBasis, New.RealPlace.differenceX₀, New.RealPlace.subframeCols, New.RealPlace.subframeBlock, New.RealPlace.volumeRatio_eq_norm_det, New.RealPlace.norm_det_subframeBlock, New.RealPlace.exists_common_vector_of_normOnePairOccursFrom, New.RealPlace.frameX₀Part, New.RealPlace.frameX₀Part_unipotentOf_inv_mul, New.RealPlace.imageBasis, New.RealPlace.dualProducts}

\begin{proof}
(a) Let $N$ be the matrix of the projection of $\gamma f$ to $yV_-$ in the basis given by the projection of $f$; then $\nu_\gamma(y)=\abs{\det N}$. Let $\omega$ be a negative basis of $yV_-\cap I_\gamma^\perp$ and $P$ the matrix of
the products of the frame with $\omega$; then $P(B_\omega)^{-1}$, with $B_\omega$ the rows of the subframe, holds the products with the basis dual to the subframe. For an index $b$ outside the subframe the
flags of $f$ and $\gamma f$ share a vector at $b$ with leading coefficients $\zeta$ and $1$, $\abs\zeta=1$, and the part of $N$ at these indices is unitriangular up to these factors; two triangular changes of basis
reduce $\abs{\det N}$ to the determinant of the products of $\omega$ with $\gamma e_1,\dots,\gamma e_i,\gamma z$ divided by that with $e_1,\dots,e_i,z$. For the lower bound take a triangular basis
(\cref{lem:triangular-basis}) and $K=\omega^*J_1\omega'$. The products of $\omega$ with the image of the subframe are $KL^{-1}B$ and those with the subframe are $L^*$, so
\[
  \nu_\gamma(y)=\abs{\det K}\,\abs\delta/\abs{\det L}^2.
\]
Expanding the form on $\omega'-\omega$, whose Gram matrix is $v^*v$ (\cref{lem:triangular-basis}(c)), gives $v^*v=\omega'^*J_1\omega'-1-K-K^*\le-1-K-K^*$, as the form is $\le0$ on the span of $\gamma f$. Thus the hermitian part of
$-K$ is at least $\frac12(1+v^*v)$, and the Ostrowski--Taussky inequality \cite{OstrowskiTaussky1951} gives $\abs{\det K}\ge\det(\frac12(1+v^*v))=2^{-i-1}\det(1+vv^*)$. If the part in $X_0$ of the image of the subframe
moves by ${}^t\bar A$ times the first $i$ rows of $B$, then $\omega'$ moves accordingly and $v$ moves by ${}^t\bar A$ times the first $i$ rows of $L$; the last column of $v$ does not move, since $L$ is lower
triangular. Left multiplication by a unipotent $u=u(a,b)$ fixing the frame adds ${}^t\bar a$ times the isotropic rows of $Z$ to the part in $X_0$ of $\gamma f$, which is a move of this kind.

(b) Fix a triangular basis at $y$ for $Z$ and write $v_0$ for the differences at $\gamma_0$, $L_1$ for the leading block of $L$, $c$ for the conjugate transpose of the first $i$ columns of $v_0$,
$G=1+v_{0,z}v_{0,z}^*$ and $C=(2^{i+1}\abs{\det L}^2/\abs\delta)^{k+i}$. Put
\[
  \phi(A)=C\,\bigl\lvert\det\bigl(G+(L_1^*A+c)^*(L_1^*A+c)\bigr)\bigr\rvert^{-(k+i)}.
\]
At a translate by $A$ the matrix $1+vv^*$ is $G+(L_1^*A+c)^*(L_1^*A+c)$, because the last column of $v$ is that of $v_0$ and the others are $v_0$'s plus ${}^t\bar A$ times the first $i$ rows of $L$; so (a)
gives $\nu_\gamma^{-(k+i)}\le\phi(A)$. Hua's bound \cite{Hua1963} for the integral of a negative power of $\abs{\det(G+Y^*Y)}$ over the rows of $Y=SA+c$, $S$ invertible and at least seven columns, gives
\[
  \int\phi\le C\,\abs{\det L_1}^{-2k}(\det G)^{-k}=\frac{(2^{i+1}\abs{\det L}^2)^i\,\bigl(2^{i+1}\abs{L_{zz}}^2/\det G\bigr)^k}{\abs\delta^{k+i}},
\]
with $\abs{\det L}=\abs{\det L_1}\abs{L_{zz}}$. Finally $\abs{\det L}\le\rho^{i+1}$ and $\abs{L_{zz}}^2/\det G\le\theta^{2(j-i)}$ (\cref{lem:triangular-basis}(b), (d)), since $\det G=1+\abs{v_{0,z}}^2$.
\end{proof}

\begin{lemma}[The lattice of columns]\label{lem:column-lattice}
For a class $c$ and a finite set $\iota$, the image of the $\iota$-tuples of vectors of $\Lambda\cap X_0$ under the conjugated embeddings is a discrete subgroup spanning $\prod_v(\C^\kappa)^\iota$ over $\R$; it has a
measurable fundamental domain of finite measure at least $1$.
\end{lemma}
\lean{New.Filling.mem_columnLattice, New.Filling.conjColumnsAtInfinity_apply, New.Filling.conjColumnsAtInfinity_injective, New.Filling.denseRange_conjColumnsAtInfinity, New.Filling.isIntegral_conjTranspose_mul, New.Filling.span_columnLattice, New.Filling.exists_int_eq_sum_re}

\begin{proof}
Hermitian products of vectors of $\Lambda\cap X_0$ are algebraic integers, as the basis of $X_0$ is orthonormal for the integral Gram matrix of the lattice. A nonzero vector has squared length an algebraic integer
whose conjugates are not all less than $1$, so the image is discrete. Every tuple of vectors of $X_0$ has a nonzero integer multiple in $\Lambda\cap X_0$ (clear the finitely many denominators), and $E$ is dense in
$\prod_v\C$, so the span is dense and closed, hence everything. The Euclidean product of two lattice points is $\sum_v\re\Phi_v(x)$ for the hermitian product $x$ of the two tuples, an algebraic integer. Let
$\iota\in E$ with $\iota^2=-1$ and let $t_v$ be the image of $x$ under the embedding at $v$ that sends $\iota$ to $\sqrt{-1}$; then $\re t_v=\re\Phi_v(x)$, and $t=\sum_vt_v$ is an algebraic integer equal to $(p-q\sqrt{-1})/2$, for the
traces $p$ of $x$ and $q$ of $\iota x$. Its squared modulus $(p^2+q^2)/4$ is an algebraic integer and rational, hence an integer, so $p$ and $q$ are even and $\sum_v\re\Phi_v(x)=p/2$ is an integer. The lattice therefore
has an integral Gram matrix in a basis, and its covolume is the square root of a positive integer.
\end{proof}

\begin{lemma}[Cauchy--Schwarz and integration in stages]\label{lem:schur-cs}
For classes $c$, $c'$ and a rational $\gamma$ write $I(c,c',\gamma,i)$ for the integral of the boxed cutoff of $c$ times the squared norm of the cut Bessel function times the Schur weight for $(c',\gamma)$ with
exponent $i$.
\begin{enumerate}[label=\textup{(\alph*)}]
\item The square of the amplitude majorant of $\gamma$ for the cut Bessel functions is at most $I(c,c',\gamma,i)\,I(c',c,\gamma^{-1},i)$.
\item $\int\beta_c\norm{W^\flat}^2\le\nu\,\#I$, with $\#I$ the number of Bessel coordinates.
\item Let $F$ be measurable and invariant under $U^0_c(M)$, let $s$ be measurable with values in the fixer of the frame, and $B$ of finite positive measure. If $\int_BF(us(x)y)\,dx\le\abs BC$ for all $u$ in the fixer
and all boxed $y$, then $\int F\cdot\text{(boxed cutoff of $c$)}\cdot\norm{W^\flat}^2\le C\nu\,\#I$.
\end{enumerate}
\end{lemma}
\lean{New.Filling.inCut_of_besselKernel_ne_zero, New.Filling.lintegral_cutoff_mul_eq_mul_right, New.Filling.lintegral_cutoff_mul_le_of_average, New.Filling.enorm_besselSystem_cut_mul}

\begin{proof}
(a) Where the integrand of the majorant,
\[
  \beta_c(y)\,\beta_{c'}(\gamma^{-1}y)\,\norm{W^\flat(y)}\,\norm{W^\flat(\gamma^{-1}y)},
\]
is nonzero, every local factor is nonzero, each vanishes outside the cut, and by the definition of
the reduced points $y$ and $\gamma^{-1}y$ are boxed, so the cutoffs equal the boxed cutoffs. Write the integrand as the geometric mean of $\beta_c(y)\norm{W^\flat(y)}^2\nu(y)^{-(k+i)}\beta_{c'}(\gamma^{-1}y)$ and
$\beta_{c'}(\gamma^{-1}y)\norm{W^\flat(\gamma^{-1}y)}^2\nu(y)^{k+i}\beta_c(y)$; since $\nu_{\gamma^{-1}}(\gamma^{-1}y)=\nu_\gamma(y)^{-1}$ (exchanging the frames inverts the ratio), the substitution $y\mapsto\gamma y$ turns the integral of the
second into $I(c',c,\gamma^{-1},i)$, and Cauchy--Schwarz gives (a).
(b) The norm of $W^\flat$ is invariant under $U_\infty$, which acts on $W^\flat$ by a unitary character. Integrate over the orbits of $U_\infty$: the integral of $\beta_c$ over an orbit is $\nu$
(\cref{prop:blocks}(a)), and $\int_{U\backslash G}\norm{W^\flat}^2\le\#I$ (\cref{thm:bessel-schur}).
(c) Unfold along the fixer $G^0$ of the frame at the infinite places: the integral is an integral over orbits $G^0a$ with a measure of total mass one on each orbit. On an orbit inside the boxed points the
Bessel norms and all factors of the boxed cutoff are constant except the cutoff $p$ for $U^0_c(M)$; for a $\Gamma$-invariant $F$ and a cutoff $p$ of a countable group $\Gamma$ acting on a unimodular group,
$\int Fp$ does not change when $F$ is translated on the right (insert the sum of the translates of $p$ and substitute in each term). Averaging over the translations $s(x)$, $x\in B$, bounds the $p$-weighted
integral of $F$ over the orbit by $C$ times that of $p$; integrating over the orbits and using (b) gives the bound.
\end{proof}

For a term $\gamma$ of cell rank $j$ over $(c,s,c')$, $i=\min(\varrho,j)$, a term $p$ of its orbit is a \emph{lattice translate} of $p_0$ with parameter $l$, a matrix with columns in $\Lambda\cap X_0$, if modulo the isotropic space
the frame of $p$ is that of $p_0$ moved in $X_0$ by $M\,l$ times the first $i$ isotropic rows of $Z(\gamma)$. The \emph{Eichler shift} with parameter $x\in\prod_v(\C^\kappa)^i$ is the unipotent element that acts on $X_0$
by the product of the Eichler maps along $e_1,\dots,e_i$ with the parameters $\overline{M/\mathrm{sc}_v(e_l)}\,x_l$, $\mathrm{sc}_v$ the scales of the normalised basis. The \emph{isotropic row index}
$\mathrm{iri}(\gamma)$ is the index, in the lattice spanned by the products of $\Lambda\cap A$ with the frame of $h_c^{-1}\gamma h_{c'}$, of the sublattice coming from the $e_a$ with $a<j$.
\lean{New.Filling.IsLatticeTranslate, New.Filling.eichlerShift, New.Filling.shiftScaling, New.Filling.kernelBound, New.Filling.initialIsotropicIndex, New.Filling.isotropicRowsIndex, New.Filling.measurable_eichlerShift, New.Filling.eichlerShift_mem, New.Filling.IsLatticeTranslate.frameX₀Part_eq, New.Filling.frameMiddle_translateByClassReps_sub, New.Filling.HasCellRank.det_subframeBlock_ne_zero, New.Filling.HasCellRank.norm_det_subframeBlock, New.Filling.HasCellRank.of_mem_doubleCoset, New.Filling.frameProducts_eq_of_mem_doubleCoset, New.Filling.termsOver.of_mem_doubleCoset, New.Filling.col_frameMiddle_sub_mem, New.Filling.col_conjTranspose_mul_mem, New.Filling.realFrameGram_ratToInf_congr, New.Filling.realFrameGram_inv_mul_ratToInf, New.Filling.norm_det_submatrix_realFrameGram, New.Filling.norm_det_shiftScaling, New.Filling.kernelBound_le, New.Filling.kernelBound_nonneg}

\begin{lemma}[The orbit as lattice translates]\label{lem:orbit-translates}
Let $\gamma$ be a term of cell rank $j\ge1$ over $(c,s,c')$ and $i=\min(\varrho,j)$.
\begin{enumerate}[label=\textup{(\alph*)}]
\item The terms of the orbit $U^0_c\gamma U^0_{c'}$ lie over $(c,s,c')$, have the cell rank $j$ and the matrix of products $Z(\gamma)$, and the middle parts of their frames differ from that of $\gamma$, row by row, by
$M$ times $\hat\OO$-combinations of the conjugates of the isotropic rows of the products of the frame of $h_c^{-1}\gamma h_{c'}$.
\item The orbit is the union of at most $\mathrm{iri}(\gamma)^k$ sets each consisting of lattice translates of one of its points.
\item At a place $v$ the local kernel of \cref{prop:volume-ratio}(b), in the parameters of the Eichler shift, has integral at most
\[
  \mathrm{kb}_v=\Bigl(\frac{(2\rho^2)^{i+1}N^j}{\abs{D_j}_v}\Bigr)^i\Bigl(\frac{(b^{13}/M^2)^iN^{(j-\varrho)_+}}{\abs{D_j}_v}\Bigr)^k,\qquad\rho=b^{37n},
\]
and $\prod_v\mathrm{kb}_v\le\tau\,(b^{13}/M^2)^{ikf}\mu_j^{-k/2}$, with $\mu_j=\mathrm{Nm}(D_j)$ divided by $N^{2f}$ at the top cell rank.
\end{enumerate}
\end{lemma}
\lean{New.Filling.HasCellRank.exists_isLatticeTranslate, New.Filling.tsum_orbit_le_of_isLatticeTranslate, New.Filling.HasCellRank.exists_majorant_lintegral_le_kernelBound, New.Filling.frameX₀Part_eichlerShift_inv_mul, New.Filling.prod_kernelBound_le, New.RealPlace.frameX₀Part_unipotentOf_inv_mul}

\begin{proof}
(a) Write $\delta=u\gamma x$ with $u\in U^0_c$, $x\in U^0_{c'}$. Moves fix the frame and keep the products, the cell rank, and the block (a move of $c$ is a term over $(c,1,c)$). The new frame is $u'=h_c^{-1}uh_c$ times the
old one in the lattice basis, and $M^{-1}(u'-1)$ is integral; off its last block of columns a row of $u'$ agrees with the identity, and the products of $e_a$ with the frame of $h_c^{-1}\gamma h_{c'}$ are the coordinates
of that frame along $e^*_a$. (b) The middle parts of two frames of the orbit differ by $k$ rows, which in a basis of $\Lambda\cap X_0$ lie in $M$ times the lattice spanned by all isotropic rows; colour each term by the
classes of these rows modulo $M$ times the sublattice of the rows $a<j$, at most $\mathrm{iri}^k$ colours. Two terms of one colour differ by $M$ times that sublattice; the rows of the leading block are independent, so the
adelic coefficients come from a rational $l$ with columns in $\Lambda\cap X_0$, and the terms are lattice translates. (c) With $\delta_u$ the leading minor in the unit frame, $\abs{\delta_u}=\abs{\kappa_v}^{2w}\abs{D_j}_v/N^j$
for $w=\sum_{a<j}(\varrho-a)$, and the shift scaling has $\abs{\det}=(M\sqrt N)^i/\abs{\kappa_v}^w$ (\cref{ss:filling-data}). Dividing the bound of \cref{prop:volume-ratio}(b) by the Jacobian $\abs{\det}^{2k}$ of the passage to the shift
parameters, the power of $\abs{\kappa_v}$ cancels in the $k$th power and may be dropped in the $i$th power ($\abs{\kappa_v}\ge1$); with $\theta=b^{6n+2}$ one has $2^{i+1}\theta^{2(j-i)}\le b^{13i}$ (as $i\ge1$, and $i=\varrho$ is
large when $j>i$). For the product, $\prod_v\abs{D_j}_v=\mathrm{Nm}(D_j)^{1/2}=N^{f(j-\varrho)_+}\mu_j^{1/2}$, and $D_j$ is a nonzero algebraic integer, so $\prod_v\abs{D_j}_v^{-i}\le1$; the remaining factor
$\bigl((2\rho^2)^{i+1}N^j\bigr)^{if}$ is at most $\tau$ because $N\le b^{60n}$.
\end{proof}

\begin{theorem}[Schur's test for one orbit]\label{thm:schur-orbit}
Let $\gamma$ be a term of cell rank $j\ge1$ over $(c,s,c')$, $i=\min(\varrho,j)$. Then:
\begin{enumerate}[label=\textup{(\alph*)}]
\item (Cauchy--Schwarz.) The square of the amplitude majorant of $\gamma$ (with the cut Bessel functions) is at most the product of the Schur integrals for $(c,c',\gamma)$ and $(c',c,\gamma^{-1})$.
\item (Orbit.) The frames of the orbit $U^0_c\gamma U^0_{c'}$ lie in at most $\mathrm{iri}^k$ cosets of a lattice, $\mathrm{iri}$ the isotropic row index of the term; summing the Schur weights over a lattice
translate and integrating in stages,
\[
  \sum_{\text{orbit}}\text{Schur integral}\le\nu\,\#I\cdot\tau^2(b^{13}/M^2)^{ikf}\bigl(\mathrm{iri}/\sqrt{\mu_j}\bigr)^k,\qquad\tau=b^{kf/4},
\]
with $\mu_j$ the normalised norm of the leading minor of order $j$.
\item Hence the square of the sum over the orbit of the absolute amplitudes is at most $(\nu\,\#I\,\tau^2(b^{13}/M^2)^{ikf})^2(N^f/\mathsf g^\sharp)^k$, with $\#I$ the number of Bessel coordinates and $\mathsf g^\sharp$ the
minor gauge ($N^f\mu_j$ divided by the two isotropic row indices).
\end{enumerate}
\end{theorem}
\lean{New.Filling.amplitudeMajorant_sq_le, New.Filling.lintegral_boxCutoff_le_of_average, New.Filling.tsum_orbit_le_of_isLatticeTranslate, New.Filling.exists_tsum_schurWeight_le, New.Filling.tsum_schurIntegral_le, New.Filling.sq_tsum_amplitudeMajorant_doubleCoset_le, New.Filling.tsum_amplitudeMajorant_orbit_le, New.Filling.HasCellRank.exists_majorant_lintegral_le_kernelBound, New.Filling.frameX₀Part_eichlerShift_inv_mul, New.Filling.eichlerShift, New.Filling.shiftScaling, New.Filling.kernelBound, New.Filling.IsLatticeTranslate, New.Filling.HasCellRank.exists_isLatticeTranslate, New.Filling.tsum_schurWeight_fibre_le, New.Filling.prod_kernelBound_le, New.Filling.tsum_schurWeight_ratToInf_mul, New.Filling.columnLattice, New.Filling.exists_isAddFundamentalDomain_columnLattice, New.Filling.integralColumns, New.Filling.conjColumnsAtInfinity, New.Filling.initialIsotropicIndex, New.Filling.lintegral_eq_lintegral_orbit, New.Filling.lintegral_classCutoff_mul_sq_le, New.Filling.div_tailScale_pow_mul_identitySecondMoment_le_cellRankSecondMoment, New.Filling.b_pow_div_M_sq_le}

\begin{proof}
(a) is \cref{lem:schur-cs}(a) with $i=\min(\varrho,j)$.

(b) Let $F$ be the sum of the Schur weights (exponent $i$) over the orbit. It is invariant under $U^0_c$, which permutes the orbit and fixes the frame, and the sum of the Schur integrals over the orbit is
$\int F\cdot(\text{boxed cutoff of }c)\cdot\norm{W^\flat}^2$. Fix $u$ in the fixer of the frame and a boxed $y$, and let $s(x)$ be the Eichler shift with parameter $x$. By \cref{lem:orbit-translates}(b) the
orbit is the union of at most $\mathrm{iri}^k$ sets $B$ of lattice translates of a point $p_0\in B$. \emph{Fibres.} Terms of $B$ with the same parameter $l$ have frames that differ only in the isotropic space, so the
points $p^{-1}y$ differ by elements of the space $V$ of $n$-tuples of isotropic vectors seen from $y$, of real dimension $2n\varrho f$; by packing (\cref{prop:packing}(b), with $R=1$ at boxed points) their weights add up
to at most $(b^{6n^2})^{\dim V}\le\tau$ times a common bound for $\prod_v\nu^{-(k+i)}$. \emph{From cosets to the lattice.} Let $\phi=\prod_v\phi_v$ be the product of the local kernels of
\cref{lem:orbit-translates}(c) for the base point $u^{-1}p_0$. By the formula for the parts in $X_0$ of the frames after an Eichler shift (\cref{prop:volume-ratio}(a)), at the point $us(x)y$ a term with parameter $l$ is a
translate of $u^{-1}p_0$ by the rescaled point $x+l$, so its ratio of volumes is bounded by $\phi(x+l)$; hence the sum of the weights over $B$ at $us(x)y$ is at most $\tau\sum_{q\in x+(\Lambda\cap X_0)^i}\phi(q)$.
\emph{From the lattice to an integral.} The mean of $x\mapsto\sum_q\phi(x+q)$ over a fundamental domain $D$ of the lattice of columns is $\int\phi/\abs D\le\prod_v\mathrm{kb}_v$, as $\abs D\ge1$ (\cref{lem:column-lattice}).
Hence the mean of $F(us(x)y)$ over $x\in D$ is at most $C=\mathrm{iri}^k\tau\prod_v\mathrm{kb}_v$, and \cref{lem:schur-cs}(c) gives the bound $C\nu\,\#I$; insert \cref{lem:orbit-translates}(c).

(c) By (a) the square of the sum over the orbit of the majorants is at most the product of the two sums of Schur integrals over the orbits of $\gamma$ and $\gamma^{-1}$; inversion maps the first orbit onto the
second, which has the same cell rank and the same $\mu_j$ (the products of $\gamma^{-1}$ are the conjugate transpose), and its isotropic row index is the other index of $\gamma$. Multiplying the two bounds of (b),
$\mathrm{iri}(\gamma)\mathrm{iri}(\gamma^{-1})/\mu_j=N^f/\mathsf g^\sharp$.
\end{proof}

\subsection{The size of a block}\label{ss:block-size}
For a term $\gamma$ of cell rank $j$ over $(c,s,c')$ put $i=\min(\varrho,j)$ and $k_\gamma=h_c^{-1}\gamma h_{c'}$. Let $\mu_j(\gamma)=\mathrm{Nm}(D_j)$, divided by $N^{2f}$ if $j=\varrho+1$; let the \emph{subframe minor} $\delta_*(\gamma)$ be the
determinant of the products of the subframe $e_1,\dots,e_j,z$ (the whole frame if $j\ge\varrho$) with its image under $\gamma$; let $\mathrm{iri}(k_\gamma)$ be the isotropic row index (\cref{ss:schur-orbit}) and
$t(k_\gamma)$, the \emph{torsion index}, the index of $\Lambda\cap A+k_\gamma(\Lambda\cap A)$ in $\Lambda\cap(A+k_\gamma A)$. The gauges of $\gamma$ are
\[
  \mathsf g^\sharp=\frac{\mathrm{Nm}(\delta_*)}{\mathrm{iri}(k_\gamma)\,\mathrm{iri}(k_\gamma^{-1})\,N^f},\qquad\mathsf g^{\max}=\max(\mathsf g^\sharp,\mathfrak g_{\rm Nm}),\qquad\mathsf g^*=M^{2if}\,\mathsf g^{\max},
\]
the minor gauge, the maximal gauge and the \emph{level gauge}. The \emph{cell second moment} $\mathcal D_j$ is $\mathcal D_{\rm main}$ at the top cell rank and $\tau^{-16}\mathcal D_{\rm id}$ below; the \emph{orbit bound} at gauge $x$ is
$\mathrm{ob}_j(x)=\tau^2(\mathcal D_j/(n_p(M^{2if}x)^k))^{1/2}$, and the \emph{size} of the orbit of $\gamma$ is $\mathrm{size}(\gamma)=\mathrm{ob}_j(\mathsf g^{\max})=\tau^2(\mathcal D_j/(n_p\mathsf g^{*k}))^{1/2}$, in units of $\nu$. A
\emph{near representative} of a set $T$ over $(c,s,c')$ is a near term of $T$ over the block that has a cell rank, is the chosen element of its orbit $U^0_c\gamma U^0_{c'}$, and lies over an agreeing $s$.
\lean{New.Filling.cellRankSecondMoment, New.Filling.orbitBound, New.Filling.termSize, New.Filling.levelGauge, New.Filling.minorGauge, New.Filling.maxGauge, New.Filling.minorNorm, New.Filling.isotropicRowsIndex, New.Filling.torsionIndex, New.Filling.latticeDisc, New.Filling.IsNearRepresentative, New.Filling.IsNearTermAt, New.Filling.IsNearRepresentativeAtCellRank, New.Filling.nearRepresentativesOver, New.Filling.cellRank, New.Filling.discriminantBound, New.Filling.secondMoment, New.Filling.linkSecondMoment, New.Filling.pairingBoundConst}

\begin{lemma}[Cell rank, minors and the subframe]\label{lem:padded}
Let $\gamma$ have cell rank $j$ and let $P_j$ be the matrix made of the first $j$ rows of $Z=Z(\gamma)$ and of rows of the identity below them.
\begin{enumerate}[label=\textup{(\alph*)}]
\item $\gamma$ has no other cell rank, $\gamma^{-1}$ has the cell rank $j$, and $\mu_j(\gamma^{-1})=\mu_j(\gamma)$.
\item $\det P_j=D_j\ne0$, and for $b\ge j$ the combination of $\gamma f$ with the coefficients of column $b$ of $P_j^{-1}$ lies in the span of $f$.
\item $\mathrm{Nm}(\delta_*)=N^{2f}\mu_j$; in particular $\delta_*\ne0$.
\item The subframe is a basis of $A^+/I_\gamma$, and its image under $\gamma$ a basis of $\gamma A^+/I_\gamma$.
\item $\mathsf g^\sharp=N^f\mu_j/(\mathrm{iri}(k_\gamma)\mathrm{iri}(k_\gamma^{-1}))$ and $\mathsf g^{\max}=\max(N^f\mu_j,\mathfrak g_{\rm Nm})/(\mathrm{iri}(k_\gamma)\mathrm{iri}(k_\gamma^{-1}))$; at the top cell rank both indices are $1$. The gauges are positive.
\end{enumerate}
\end{lemma}
\lean{New.Filling.HasCellRank.inv, New.Filling.minorNorm_inv, New.Filling.HasCellRank.paddedLeading_mul_inv, New.Filling.HasCellRank.abs_norm_det_updateRow_paddedLeading, New.det_updateRow_of_mulVec_eq_single, New.Filling.isotropicRowsIndex_top, New.Filling.isotropicRows_conjTranspose_mul_eq, New.Filling.HasCellRank.cellRank_eq, New.Filling.HasCellRank.hasCellRank_cellRank, New.Filling.HasCellRank.frameProducts_mulVec_eq_zero_of_shared, New.Filling.HasCellRank.eq_zero_of_frameProducts_mulVec_eq_zero, New.Filling.HasCellRank.eq_zero_of_mem_commonIsotropic, New.Filling.HasCellRank.exists_mem_commonIsotropic, New.Filling.frameImage_one_mulVec_mem_commonIsotropic_iff, New.linearIndependent_mkQ_col_and_colSpan_sup_eq, New.mulVec_eq_submatrix_mulVec_of_eq_zero_off, New.Filling.isIntMatrix_conjTranspose_mul_of_classCoordX₀, New.Filling.one_le_latticeDisc, New.Filling.card_eq_finrank_colSpan, New.Filling.classCoord_eq_fromBlocks, New.Filling.classCoord_gram, New.Filling.classCoord_mul_framePair, New.Filling.classCoord_mul_isotropicFramePair, New.Filling.framePair_translateByClassReps_mul_classCoord, New.Filling.gram_classCoordX₀, New.Filling.isIntMatrix_framePair}

\begin{proof}
(a) Let $j<j'$ both be cell ranks. Then $j'\le\varrho$, since the top rank is decided by $\mathfrak g\ne0$. The shared vector at the index $j'-1\ge j$ is $fx=\gamma fy$ with $y$ supported on the first $j'$ indices and
$y_{j'-1}=1$; its products with $e_a$, $a<j'$, vanish because $A$ is totally isotropic and orthogonal to $z$, so the leading block of order $j'$ of $Z$ kills the nonzero vector $y$, against $D_{j'}\ne0$. The
products of $\gamma^{-1}$ form $Z(\gamma)^*$, so its leading minors are the conjugates, $\mathfrak g$ is unchanged, and a shared vector $fx=\gamma fy$ gives $\gamma^{-1}fx'=fy'$ after exchanging the roles of the two frames.
(b) Expanding along the rows of the identity leaves the leading minor. For $b\ge j$, combining the shared vectors at the indices $j,\dots,\varrho$ (each with coefficient $1$ at its top index) gives $y$ with
$y_l=\delta_{lb}$ for $l\ge j$ and $\gamma fy$ in the span of $f$; as in (a) the first $j$ entries of $Zy$ vanish, so $P_jy$ is the $b$th unit vector.
(c) At the top rank $\delta_*=D_{\varrho+1}$. Below it, $\delta_*$ is the determinant of $P_j$ with its last row replaced by the last row of $Z$. The last column $y$ of $P_j^{-1}$ gives a shared vector
$\gamma z+\dots=\zeta z+\dots$ with $\zeta\bar\zeta=1$, and the last row of $Z$ sends $y$ to $h(z,\zeta z)=-N\zeta$; replacing a row of a matrix whose inverse has $y$ as that column multiplies the determinant by the new row
applied to $y$. Hence $\delta_*=-N\zeta D_j$ and $\mathrm{Nm}(\delta_*)=N^{2f}\mathrm{Nm}(D_j)$.
(d) Let $\gamma fx'\in I_\gamma$ with $x'$ supported on the subframe. It lies in $\gamma A$, and the columns of $\gamma f$ are independent, so $x'_z=0$; it lies in $A$, so its products with $e_a$, $a<j$, which are the
entries $(Zx')_a$, vanish; as $D_j\ne0$, $x'=0$. Coefficients outside the subframe are absorbed by $I_\gamma$: if $\mathfrak g=0$, $A=A_{\le j}+I_\gamma$ and $\gamma A=\gamma A_{\le j}+I_\gamma$ (\cref{prop:links}(c)), and if
$\mathfrak g\ne0$ there are none. The statement for $A^+$ follows by applying this to $\gamma^{-1}$, which carries $I_{\gamma^{-1}}$ to $I_\gamma$.
(e) is (c) and the definitions; at the top rank all isotropic vectors are counted, and $\mathfrak g\ne0$ there. The indices are finite, hence positive (\cref{thm:link-lattice} below), and $\mu_j>0$ by (b).
\end{proof}

\begin{lemma}[Orbit bounds]\label{lem:orbit-bound}
\begin{enumerate}[label=\textup{(\alph*)}]
\item $\abs{\mathfrak g(\gamma)}_v\le\abs{D_{\varrho+1}}_v^2$, and for a near term $\mathfrak g_{\rm Nm}\le b^{270n^2f}$ and $\mathsf g^*\le b^{n^3f}$.
\item $\mathcal D_j\le\mathcal D_{\rm main}$, $n_p\le\mathcal D_{\rm main}$, and $\mathcal D_j\ge\tau^{8-32i}\mathcal D_{\rm id}$ for $j\ge1$.
\item If $\gamma$ is near and $\mathfrak g_{\rm Nm}\ge1$, then $\mathrm{ob}_{\varrho+1}(\mathfrak g_{\rm Nm})\ge b^{-k^3f}$, and the product over the infinite places of the local bounds of \cref{thm:open-cell}(b) at $\gamma$ is at
most $(\nu\,\mathrm{ob}_{\varrho+1}(\mathfrak g_{\rm Nm})/2)^2$.
\item If $\gamma$ has cell rank $j\ge1$, lies over $(c,s,c')$ and $\mathfrak g_{\rm Nm}\le N^f\mu_j$, then the absolute amplitudes of its orbit add up to at most $\nu\,\mathrm{ob}_j(\mathsf g^{\max})$.
\end{enumerate}
\end{lemma}
\lean{New.Filling.norm_map_termDisc_le, New.Filling.IsNear.termDiscNorm_le, New.Filling.IsNear.levelGauge_le, New.Filling.cellRankSecondMoment_nonneg, New.Filling.orbitBound_top, New.Filling.pairCountScale_le_mainSecondMoment, New.Filling.div_tailScale_pow_mul_identitySecondMoment_le_cellRankSecondMoment, New.Filling.cellRankSecondMoment_le_mainSecondMoment, New.Filling.IsNear.inv_pow_le_orbitBound, New.Filling.pairingBoundConst, New.Filling.prod_pairingBoundSq_eq, New.Filling.prod_frameScaleSq_bounds, New.Filling.pairingBoundConst_le_mul, New.Filling.two_pow_mul_pairCountScale_mul_le, New.Filling.pairingBoundConst_le, New.Filling.prod_pairingBoundSq_le, New.Filling.b_pow_div_M_sq_le, New.Filling.tsum_amplitudeMajorant_orbit_le}

\begin{proof}
(a) $\mathfrak g=\abs{D_{\varrho+1}}^2-N^2\abs{D_\varrho}^2$ is a difference of two nonnegative numbers at each place, and it is $\ge0$ (\cref{prop:links}(b)). For a near term the entries of $Z$ are at most $b^{134n}$
(\cref{prop:links}(b)), so by Hadamard's inequality the squares of its minors are at most $b^{270n^2}$ at every place; the indices in the gauges are at least $1$, and $M^{2if}N^f\le b^{(46n+60n)f}$, with
$46n+60n+270n^2\le n^3$. (b) $\tau^{-16}\mathcal D_{\rm id}\le\mathcal D_{\rm id}\le\tau^{44}\mathcal D_{\rm id}\le\mathcal D_{\rm main}$ by admissibility, $n_p\le\mathcal D_{\rm id}$ as $\#I\ge1$, and for $i\ge1$, $\tau^{8-32i}\le\tau^{-24}\le\tau^{-16}$,
while $\mathcal D_{\rm main}\ge\mathcal D_{\rm id}$.
(c) Comparing squares, the pair count $n_p$ cancels against $n_p\le\mathcal D_{\rm main}$, and it remains that $M^{2\varrho fk}\mathfrak g_{\rm Nm}^k\le b^{(23\cdot2\varrho+270n^2)fk}\le b^{2k^3f}$, since $\varrho+1\le n\le k$. For the
local bounds: their product is a constant $C_{\rm pb}$ divided by $\mathfrak g_{\rm Nm}^{10^5n^4+k-1}\ge\mathfrak g_{\rm Nm}^k$. At a place the squared scales $s_a=\kappa_v^{2(\varrho-a)}/N$ of the frame have product between $b^{-60n^2}$ and
$b^{46n^2}$, and $\rho^{9n}/s^2\le b^{453n^2}$ for $\rho=b^{37n}$; with $(453\cdot10^5+92+184)n^6f\le10^5n^9f$ this gives
$C_{\rm pb}\le2^{kf}\tau(\#I)^2/(\deg\,c_G^f\Lambda^2)$, and with the bound for the scaled Siegel constant of admissibility (\cref{prop:filling-admissible}) and $2^{kf}n_pM^{2\varrho fk}\le\tau a_*\vol(\hat\OO)^n$, one gets
$C_{\rm pb}\le(\nu\tau^2/2)^2\mathcal D_{\rm main}/(n_pM^{2\varrho fk})$; dividing by $\mathfrak g_{\rm Nm}^k$ gives the claim.
(d) By \cref{thm:schur-orbit}(c), and as $\mathsf g^{\max}=\mathsf g^\sharp$ under the hypothesis, it suffices that $(\#I)^2(b^{13}/M^2)^{2ikf}N^{kf}\le\mathcal D_j/(n_pM^{2ikf})$. Since $b^{19}\le M$, $b^{13}/M^2\le(b^4M)^{-1}$, so the left side
is at most $(\#I)^2N^{kf}b^{-8ikf}M^{-2ikf}$. By (b), $\mathcal D_j/n_p\ge b^{2kf-8ikf}\mathcal D_{\rm id}/n_p=b^{2kf-8ikf}\#I\,N^{kf}$, as $\mathcal D_{\rm id}=\#\bar R_0n_X\#I$; and $\#I\le b^{2kf}$ by admissibility.
\end{proof}

\begin{theorem}[Size of a block]\label{thm:block-size}
\begin{enumerate}[label=\textup{(\alph*)}]
\item The second moment of the difference between the blocks of $U$ and those for the cut functions is at most $b^{-40}\mathcal D_{\rm id}$; $\sqrt{\text{second moment}}$ is subadditive.
\item (Open cell.) For a near term with $\mathfrak g\ne0$ and a measurable weight $\omega\le1$ invariant under the fixer of the frame and zero at the points in the cut but not in the box, the integral of $\omega$
times the cut kernel has norm at most the root of the product of the local bounds of \cref{thm:open-cell}(b), plus the integral of $\omega$ times the difference of the kernels; and this is at most the
orbit bound at the top cell rank.
\item For a stable $T\subseteq\Gamma^*$, the cut block of $T$ over $(c,s,c')$ is at most $\nu$ times the sum of the sizes of the near representatives over the block (the near terms that are the chosen
element of their orbit and have a cell rank); there are none unless $s$ agrees.
\end{enumerate}
\end{theorem}
\lean{New.Filling.secondMoment_sub_cut_le, New.Filling.sqrt_secondMoment_add_le, New.Filling.secondMoment_nonneg, New.Filling.sum_norm_amplitudeSum_sub_cut_le, New.Filling.card_residueGL_mul_b_pow_le, New.Filling.enorm_integral_besselKernel_le, New.Filling.enorm_amplitudeSum_le_orbitBound_top, New.Filling.enorm_amplitudeSum_le_orbitBound, New.Filling.enorm_amplitudeSum_le_tsum_termSize, New.Filling.enorm_amplitudeSum_le_of_termDiscNorm_le, New.Filling.hasSum_amplitudeSum_fibre, New.Filling.enorm_amplitudeSum_le_tsum_fibre, New.Filling.hasCellRank_cellRank_of_mem, New.Filling.isNear_of_amplitudeMajorant_ne_zero, New.Filling.exists_isNearRepresentative_of_besselKernel_ne_zero, New.Filling.tsum_ite_representative_eq_one, New.Filling.tsum_representative_weight_eq_one, New.Filling.isIntegral_termDisc, New.Filling.eq_one_of_mul_eq_mul_of_termDisc_ne_zero, New.Filling.toReal_smul_besselKernel_eq_zero, New.Filling.amplitudeSum_eq_inter_compl_add_inter, New.Filling.IsNear.inv_pow_le_orbitBound, New.Filling.prod_pairingBoundSq_le, New.Filling.prod_pairingBoundSq_eq, New.Filling.prod_frameScaleSq_bounds, New.Filling.pairingBoundConst_le_mul, New.Filling.two_pow_mul_pairCountScale_mul_le, New.Filling.pairingBoundConst_le, New.Filling.norm_map_termDisc_le, New.Filling.IsNear.termDiscNorm_le, New.Filling.cellRankSecondMoment_nonneg, New.Filling.orbitBound_top, New.Filling.pairCountScale_le_mainSecondMoment}

\begin{proof}
(a) For fixed $c$ and $s$ the norms of the differences of the blocks, summed over $c'$, are at most $\eta=\nu b^{-k^3f}/2$ (\cref{prop:blocks}(b)); so the sum of their squares is at most $\eta^2$. Summing over $s\in\bar G$ and over the
$n_X\le\mathcal D_{\rm id}$ classes $c$ and dividing by $\nu^2$ gives at most $\#\bar G\,b^{-2k^3f}\mathcal D_{\rm id}/4$. The group $\bar G$ is a product of $f$ general linear groups of degree $2\varrho+k+1$ over each of the two residue fields
of order at most $p_i$, so $\#\bar G\le M^{(2\varrho+k+1)^2f}\le b^{23(2\varrho+k+1)^2f}$, and $23(2\varrho+k+1)^2f+40\le2k^3f$. The square root of the second moment is a multiple of the norm in an $\ell^2$ space of blocks, hence
subadditive.

(b) The first statement is \cref{thm:open-cell}(c) for the weight $\omega$. As $\mathfrak g$ is a nonzero algebraic integer, $\mathfrak g_{\rm Nm}\ge1$; by \cref{lem:orbit-bound}(c) the root of the product of the local bounds is
at most $\nu\,\mathrm{ob}_{\varrho+1}(\mathfrak g_{\rm Nm})/2$, and the integral of $\omega$ times the difference of the kernels is at most $\nu b^{-k^3f}/2\le\nu\,\mathrm{ob}_{\varrho+1}(\mathfrak g_{\rm Nm})/2$.

(c) \emph{Terms with $\mathfrak g\ne0$ throughout.} If $p\gamma=\gamma w$ with $p\in U^0_c$, $w\in U_{c'}$ and $\mathfrak g(\gamma)\ne0$, then $w$ fixes the frame of $\gamma^{-1}$, whose matrix of products with $f$ is invertible; so $w$ fixes vectors
whose dual coordinates form the identity matrix, $w=1$ and $p=1$. Hence in each double coset $U^0_c\gamma U_{c'}$ exactly one pair $(a,b)$ carries $\gamma$ to the chosen representative $a\gamma b^{-1}$, where the representative is chosen
to be a near representative if the double coset contains one. The weight $\phi_\gamma(y)=[\gamma\text{ is the representative}]\,\omega(y)$, $\omega$ the indicator of the reduced points, then has orbit sums $1$ (the reduced
indicator is a partition of unity in boxes, \cref{prop:packing}(c)), and \cref{prop:blocks}(b) lets it replace $\beta_c(y)\beta_{c'}(\gamma^{-1}y)$. If the representative $\gamma$ is a near representative, (b) bounds its integral by $\nu\,\mathrm{ob}_{\varrho+1}(\mathfrak g_{\rm Nm})$. Otherwise the cut
kernel of $\gamma$ vanishes at every reduced point: if it did not vanish at a reduced $y$, translating $\gamma^{-1}y$ into the boxed points by some $u\in U_{c'}$ (\cref{prop:packing}(c)) and taking the chosen element of the
orbit of $\gamma u^{-1}$ would give a near representative in the double coset.

\emph{Terms of one cell rank.} Let $T$ be stable with all elements of cell rank $j$ and none of cell rank $0$ over the block. On the stable part $K$ where $\mathfrak g_{\rm Nm}>N^f\mu_j$ the cell rank is the top one,
the indices are $1$ and $\mathsf g^{\max}=\mathfrak g_{\rm Nm}$, so the previous paragraph applies with $\mathrm{ob}_{\varrho+1}(\mathfrak g_{\rm Nm})=\mathrm{ob}_j(\mathsf g^{\max})$. Off $K$, sum the absolute amplitudes orbit by orbit: an orbit
whose sum is nonzero contains a term with a nonzero majorant, which is near (otherwise the boxed cutoffs at $y$ and $\gamma^{-1}y$ are never both nonzero and the Schur integrals vanish), and its chosen element is a near
representative, whose orbit contributes at most $\nu\,\mathrm{ob}_j(\mathsf g^{\max})$ by \cref{lem:orbit-bound}(d). The blocks vanish unless $s$ agrees (\cref{prop:blocks}(c)).

\emph{Sorting by the cell rank.} Every element of $\Gamma^*$ has a cell rank (\cref{lem:common-vectors}(c) for $\mathfrak g=0$; the top rank for $\mathfrak g\ne0$), the parts of $T$ of one cell rank are stable (\cref{prop:links}(a)), and
none has cell rank $0$, since a norm-one pair from $0$ would put $\gamma$ off $\Gamma^*$. The norm of the block is at most the sum of the norms of the blocks of these parts.
\end{proof}

\subsection{The lattice of a link and the gain}\label{ss:gain}
For a term $\gamma$ over $(c,s,c')$ let $\Lambda$ be the lattice of $c$. The vectors of $\Lambda$ orthogonal to $A$ and $\gamma A$ give a lattice $L^1_\gamma$ in the link; those whose $A$-part is divisible by $M$ give
$L^+_\gamma$; and the \emph{lattice of the term} $L_\gamma\subseteq\Lambda\cap X_0$ consists of the $v$ for which some vector of $v+M(\Lambda\cap A)$ is orthogonal to $\gamma f$. For two terms, their \emph{common link} is the
subspace of $X_0$ sent by both links to one vector of $X_0$, and their \emph{common lattice} $L_{12}$ is the part of $L_\gamma\cap L_{\gamma'}$ in it; its corank $l$ is $k$ minus its rank.
\lean{New.Filling.classLatticeX₀, New.Filling.termLattice, New.Filling.commonPart, New.Filling.pairLattice, New.Filling.pairDiscriminant, New.Filling.linkCorank, New.Filling.HasIntegralLifts, New.Filling.HasLevelLifts, New.Filling.linkMetric, New.Filling.linkDiscSq, New.Filling.linkMetric_posDef, New.Filling.LinkSplitting, New.Filling.linkVectors, New.Filling.linkVectorPairs, New.Filling.IsOrthogonalLift, New.Filling.subframeMinor, New.Filling.subframeIndices, New.Filling.paddedLeading, New.Filling.det_paddedLeading, New.Filling.isotropicRows, New.Filling.frameMiddle, New.Filling.IndexBelow, New.Filling.IndexNotBelow, New.submoduleIndex, New.colLattice, New.genIndex}

On pairs of middle coordinates $X_0\oplus Ez$ write $T=\mathrm{diag}(1,-N)$ for the middle Gram matrix, $\mathrm{lf}=\mathrm{diag}(T,-T)$ for the \emph{link form} and
$Q=\frac12\mathrm{diag}(T,T)+N^{-1}L^*L$, with $L$ the two $z$-rows of $T$, for the \emph{link metric}. For a matrix $v$ with columns in the link and $e\in\GL(\Ad_f)$ write $\mathrm{lds}(v,e)=\mathrm{Nm}(\det v^*Qv)[\det e]^2$, the
squared discriminant of the lattice with basis $ve$ for $Q$. The \emph{$z$-forms} of a pair are the products of its two members with $z$.
\lean{New.Filling.zFormIndex, New.Filling.zFormIndex_injective, New.Filling.lastFramePair, New.Filling.lastPair, New.Filling.listPair, New.Filling.eStarCols, New.Filling.classOf, New.Filling.sourceCols}

\begin{lemma}[Link forms and Liouville]\label{lem:liouville}
\begin{enumerate}[label=\textup{(\alph*)}]
\item For $t$ orthogonal to $A$, $\gamma A$ and $t'$ orthogonal to $A$, $\gamma'A$, the pairs $p$, $p'$ of $t$ and $t'$ in the links of $\gamma$ and $\gamma'$ have $\mathrm{lf}(p,p')=h(t,t')-h(t,\gamma\gamma'^{-1}t')$; in particular the link is
isotropic for the link form.
\item $Q$ is the mean of the forms of $X_0\perp Ez^+$ on the two images $\mathrm{src}$ and $\mathrm{tgt}$ of a pair; on the link it is the form of $X_0\perp Ez^+$ on the source; it is positive definite at every infinite
place, and independent columns have a nonzero Gram determinant.
\item (Liouville.) Let $Q'$ be positive definite at every infinite place, $a$, $b$ matrices over $E$ with $\det(a^*Q'b)\ne0$, $v$, $v'$ invertible finite adelic matrices, and $0\le x\le[\det(v^*(a^*Q'b)v')]$. Then
\[
  x^2\le[\det v]^2\,\mathrm{Nm}(\det a^*Q'a)\,\mathrm{Nm}(\det b^*Q'b)\,[\det v']^2.
\]
\end{enumerate}
\end{lemma}
\lean{New.Filling.liouville_gram, New.Filling.gram_linkMatrix_mul, New.Filling.framePair_last_eq_zRowBlockDiagonal_mul, New.Filling.dotProduct_linkMetric_mulVec, New.Filling.linkMetric_form_of_mem_link, New.Filling.linkMetric_eq_smul_fromBlocks, New.Filling.det_gram_linkMetric_ne_zero, New.Filling.positiveGram_posDef, New.Filling.gram_finAdeleMatrix_mul, New.Filling.conjTranspose_mul_mul_apply_eq_dotProduct, New.Filling.gram_linkMetric_eq_gram_sourceCols, New.Filling.gram_linkForm_eq_gram_sourceCols_sub}

\begin{proof}
(a) The columns of $t$ and $\gamma^{-1}t$ have vanishing coordinates at the $e^*_a$, and on such vectors the form is $T$ on the middle coordinates; so $\mathrm{lf}(p,p')=h(t,t')-h(\gamma^{-1}t,\gamma'^{-1}t')$, and $h(\gamma^{-1}t,\gamma'^{-1}t')=h(t,\gamma\gamma'^{-1}t')$
since $\gamma$ is an isometry. For $\gamma'=\gamma$ the two terms cancel. (b) Expand: $\frac12(T\oplus T)$ gives the mean of the forms of the two members, and $N^{-1}\abs{\text{$z$-form}}^2$ on each half turns $-N\abs\lambda^2$ into
$+N\abs\lambda^2$ after the exchange of the $z$-coordinates that defines $\mathrm{src}$ and $\mathrm{tgt}$. On the link $\mathrm{src}$ and $\mathrm{tgt}$ differ by the isometry $\Upsilon_\gamma$ (\cref{lem:link-isometry}). In the coordinates $Q$ is
diagonal with entries $1/2$ on $X_0$ and $N/2$ at $z$. (c) At each place the Gram matrix of the columns of $(a\mid b)$ for $Q'$ is positive semidefinite, which gives
$\abs{\det a^*Q'b}_v^2\le\det(a^*Q'a)_v\det(b^*Q'b)_v$; multiplying over the places and squaring, $\mathrm{Nm}(\det a^*Q'b)^2\le\mathrm{Nm}(\det a^*Q'a)\mathrm{Nm}(\det b^*Q'b)$. Since $[\det(v^*Xv')]=[\det v]\mathrm{Nm}(\det X)[\det v']$ for $X$ over $E$,
the hypothesis gives $x\le[\det v]\mathrm{Nm}(\det a^*Q'b)[\det v']$.
\end{proof}

\begin{lemma}[The vectors $z^+$ and $\Upsilon_\gamma^{-1}z^+$]\label{lem:link-vectors}
Let $\gamma$ be a term over $(c,s,c')$ and $X_0^\gamma=X_0\cap\Upsilon_\gamma^{-1}X_0$.
\begin{enumerate}[label=\textup{(\alph*)}]
\item The products of $a\in X_0\perp Ez^+$ with $z^+$ and $\Upsilon_\gamma^{-1}z^+$ are $N$ times the $z^+$-coordinates of $a$ and $\Upsilon_\gamma a$; the vectors orthogonal to both are those of $X_0^\gamma$. For $y$ orthogonal to $A$ and $\gamma A$
with source $a$, these products are minus the products of $y$ with $\gamma z$ and $z$.
\item The products of $z^+$ and $\Upsilon_\gamma^{-1}z^+$ with the images of vectors of $\Lambda$ orthogonal to $A$ and $\gamma A$ are integral.
\item A vector $v\in X_0$ lies in $L^+_\gamma\cap X_0^\gamma$ if and only if $v\in\Lambda$ and some vector of $v+M(\Lambda\cap A)$ is orthogonal to $\gamma f$; hence $L_\gamma=L^+_\gamma\cap X_0^\gamma$.
\item Let $\gamma$ have cell rank $j$ and let $u$ be the $z^+$-coordinate of $\Upsilon_\gamma^{-1}z^+$. Then $\mathfrak g=\abs{D_j}^2(1-\abs u^2)$; the Gram matrix of $z^+$, $\Upsilon_\gamma^{-1}z^+$ is $N\bigl(\begin{smallmatrix}1&u\\\bar u&1\end{smallmatrix}\bigr)$, with
determinant $N^2\mathfrak g/\abs{D_j}^2$. If $\mathfrak g=0$, $\Upsilon_\gamma^{-1}z^+$ is a multiple of $z^+$; if $\mathfrak g\ne0$, the two are independent and the product over the places of the moduli of their Gram
determinant is $N^f\mathfrak g_{\rm Nm}/\mathsf g^\sharp$.
\item The pairs $\varepsilon$ of the link with sources $-\Upsilon_\gamma^{-1}z^+$ and $-z^+$ represent the $z$-forms: $\varepsilon^*Qw$ is the matrix of the $z$-forms of the columns $w$ of the link, and the Gram matrix of $\varepsilon$ for
$Q$ is that of $\Upsilon_\gamma^{-1}z^+$, $z^+$ in $X_0\perp Ez^+$.
\end{enumerate}
\end{lemma}
\lean{New.Filling.linkVectors_mulVec, New.Filling.frameImage_conjTranspose_mulVec_eq_zero_iff, New.Filling.linkIsometry_perpToPositive_apply, New.Filling.linkVectors_mulVec_perpToPositive, New.Filling.isIntMatrix_lastFramePair_mul, New.Filling.elim_mem_preimageLattice_classCoord_iff, New.Filling.HasCellRank.termDisc_eq_mul_linkIsometry_symm, New.Filling.gram_linkVectors, New.Filling.HasCellRank.det_gram_linkVectors_mul, New.Filling.HasCellRank.norm_det_gram_linkVectors, New.Filling.prod_norm_det_gram_linkVectors, New.Filling.linkVectorPairs, New.Filling.col_linkVectorPairs_mem, New.Filling.linkSource_col_linkVectorPairs, New.Filling.conjTranspose_linkVectorPairs_mul_mulVec_apply, New.Filling.col_mem_of_colSpan_le, New.Filling.linkVectorPairs_mul_linkMetric_mul, New.Filling.gram_linkVectorPairs, New.Filling.linkVectors}

\begin{proof}
(a) The first statement is the definition of the positive form, in which $z^+$ has norm $N$, and of $\Upsilon_\gamma$, an isometry. For $y$ in $(A+\gamma A)^\perp$ with pair $([y],[\gamma^{-1}y])$, the $z^+$-coordinate of the source
is the $z$-coordinate of $\gamma^{-1}y$ and that of the target, $\Upsilon_\gamma a$, is the $z$-coordinate of $y$; the products of $y$ with $z$ and of $\gamma^{-1}y$ with $z$ are $-N$ times these. (b) In the lattice basis, where the middle
block of the Gram matrix is integral, the linear forms $x\mapsto h(x,k_\gamma z)$ and $x\mapsto h(x,z)$ have integral coefficients when $k_\gamma$ stabilises the standard lattice. (c) Write the vector as $a+v$ with $a\in A$ and
$v\in X_0$; it lies in $\Lambda$ exactly when $a$ has integral coordinates and $v\in\Lambda\cap X_0$. A vector of $X_0$ in $X_0^\gamma$ has vanishing $z$-forms, and a lift $y=v-Ma$ orthogonal to $A$ and $\gamma A$ with
vanishing $z$-coordinate in $\gamma^{-1}y$ is orthogonal to $\gamma f$. (d) If $\mathfrak g=0$, then $X_0^\gamma=X_0$ (\cref{prop:links}(c)), so $\Upsilon_\gamma^{-1}z^+$ is orthogonal to $X_0$, hence $\mu z^+$, and $\abs\mu=1$ because $\Upsilon_\gamma$ is an isometry; both sides of the identity
vanish. If $\mathfrak g\ne0$, then $j=\varrho+1$ and $D=\det Z\ne0$; Cramer's rule gives $y$ with $Zy=De_z$ and $y_z=D_\varrho$, and the combination $\gamma fy$, which is orthogonal to the $e_a$ and has product $D$ with
$z$, shows that the link sends $y_0+D_\varrho z^+$, for some $y_0\in X_0$, to $-(D/N)z^+$. Hence $u=-ND_\varrho/D$ and $\abs D^2(1-\abs u^2)=\abs D^2-N^2\abs{D_\varrho}^2=\mathfrak g$. The Gram matrix follows from (a), and its determinant is
$N^2(1-\abs u^2)$. If $\mathfrak g\ne0$, $\abs u\ne1$ and the vectors are independent; at the top cell rank
$\mathsf g^\sharp=\mathrm{Nm}(D_{\varrho+1})/N^f$, and the product of $N^2\abs{\mathfrak g}_v/\abs{D_{\varrho+1}}_v^2$ over the places is $N^{2f}\mathfrak g_{\rm Nm}/\mathrm{Nm}(D_{\varrho+1})=N^f\mathfrak g_{\rm Nm}/\mathsf g^\sharp$. (e) By (b) of \cref{lem:liouville}
on the link $Q$ is the positive form on the sources, and by (a) the products of the source of a pair with $-\Upsilon_\gamma^{-1}z^+$ and $-z^+$ are its two $z$-forms.
\end{proof}

An \emph{orthogonal lift} of a basis $v$ of the link of $\gamma$ at cell rank $j$ is a matrix $u$ whose columns are orthogonal to $A$ and $\gamma A$, whose first columns, one for each $b\ge j$, are dual to the
$e^*_b$ and sent to $0$ in the link (they are a basis of $I_\gamma$), and whose other columns are orthogonal to the $e^*_b$ and sent to the columns of $v$.
\lean{New.Filling.IsOrthogonalLift, New.Filling.HasCellRank.exists_isotropic_inter_vector, New.Filling.IsOrthogonalLift.apply_inl, New.Filling.IsOrthogonalLift.conjTranspose_isotropicRows_mul_eq_zero, New.Filling.IsOrthogonalLift.dual, New.Filling.IsOrthogonalLift.gram_add_border, New.Filling.IsOrthogonalLift.linkMatrix_mul, New.Filling.IsOrthogonalLift.linkMatrix_mul_toCols₁, New.Filling.IsOrthogonalLift.orthogonal, New.Filling.IsOrthogonalLift.orthogonal_toCols₁, New.Filling.IsOrthogonalLift.toCols₁_eq_fromRows, New.Filling.IsOrthogonalLift.toCols₁_mul_framePair, New.Filling.IsOrthogonalLift.toCols₁_mul_self, New.Filling.toRows₁_gram_framePair_isotropicFramePair, New.Filling.toRows₂_gram_framePair_isotropicFramePair, New.Filling.submoduleIndex_colLattice_dvd_preimageLattice, New.Filling.IsOrthogonalLift.exists_basis_orthogonal, New.Filling.classCoord_orthogonal, New.Filling.gram_classOf, New.Filling.IsOrthogonalLift.toCols₁_mul_toCols₂, New.Filling.IsOrthogonalLift.toCols₁_mul_isotropicFramePair, New.Filling.IsOrthogonalLift.toCols₁_mul_lastPair, New.Filling.classCoord_mul_listPair, New.Filling.IsOrthogonalLift.exists_mul_classOf_eq_and_linkMatrix_mul_eq, New.Filling.IsOrthogonalLift.col_mem_colLattice_isotropicFramePair, New.Filling.IsOrthogonalLift.gram_classOf_add_border, New.Filling.IsOrthogonalLift.eq_mul_of_isotropicRows_mul_eq_zero}

\begin{theorem}[The lattice of a link]\label{thm:link-lattice}
Let $\gamma$ be a term of cell rank $j$ over $(c,s,c')$, $i=\min(\varrho,j)$.
\begin{enumerate}[label=\textup{(\alph*)}]
\item $\gamma$ has orthogonal lifts of a basis of its link, isometric to it for the link metric after a correction by the products with $z$ and $\gamma z$; adjoining two vectors of norm $N$ gives isotropic
lists $(e_a,z+z')$, $(\gamma e_a,\gamma z+z'')$, $a<j$, whose Gram determinant has norm the square of the norm of the subframe minor $\delta_*(\gamma)\ne0$ (Jacobi's identity in the quotient).
\item There is a basis of the link in $L^+_\gamma$ whose torsion index times discriminant is at most $M^{2if}\mathsf g^\sharp/N^f$.
\item The products of $z^+$ and $\Upsilon_\gamma^{-1}z^+$ with $L^1_\gamma$ are integral; their orthogonal complement is $X_0\cap\Upsilon_\gamma^{-1}X_0$; they are proportional if $\mathfrak g=0$ and independent otherwise,
with an explicit Gram determinant.
\end{enumerate}
Moreover the indices $\mathrm{iri}(k_\gamma)$, $\mathrm{iri}(k_\gamma^{-1})$ and $t(k_\gamma)$ are finite, and $\mathrm{iri}(k_\gamma)\mathrm{iri}(k_\gamma^{-1})t(k_\gamma)$ divides the index of the integral classes of the two lists.
\end{theorem}
\lean{New.Filling.exists_orthogonal_lift, New.Filling.IsOrthogonalLift.linkDiscSq_mul_sq_eq, New.Filling.exists_lifts_linkDiscSq_le, New.Filling.exists_hasLevelLifts_linkDiscSq_le, New.Filling.abs_norm_det_gram_eq_abs_norm_subframeMinor_sq, New.Filling.isotropicRowsIndex_mul_dvd_submoduleIndex, New.Filling.torsionIndex_dvd_submoduleIndex, New.Filling.HasCellRank.subframeMinor_ne_zero, New.Filling.HasCellRank.linearIndependent_mkQ_frameImage_and_sup_eq, New.Filling.HasCellRank.linearIndependent_mkQ_frameImage_one_and_sup_eq, New.Filling.HasCellRank.linearIndependent_mkQ_frameImage_coe_and_sup_eq, New.Filling.linkVectors_mulVec_perpToPositive_mem, New.Filling.linkVectors_mulVec_eq_zero_iff, New.Filling.exists_lift_mem_linkedX₀_iff_exists_orthogonal, New.Filling.HasCellRank.exists_linkIsometry_symm_eq_smul, New.Filling.linearIndependent_linkVectors_and_prod_norm_det, New.Filling.linearIndependent_linkVectors, New.Filling.prod_norm_det_gram_linkVectors, New.Filling.linkVectorPairs_mul_linkMetric_mul, New.Filling.gram_linkForm_eq_zero_of_le_link, New.Filling.HasCellRank.exists_shared_frame_vector, New.Filling.HasCellRank.minorGauge_eq, New.Filling.HasCellRank.maxGauge_eq, New.Filling.HasCellRank.abs_norm_subframeMinor, New.Filling.HasCellRank.minorNorm_pos, New.Filling.HasCellRank.unique, New.Filling.levelGauge_pos, New.Filling.isotropicRowsIndex_pos, New.Filling.torsionIndex_pos, New.Filling.maxGauge_pos}

\begin{proof}
(a) \emph{Lifts.} The three conditions on a column of $u$ form one linear system. For a column of $(0\mid v)$ in the link there is $t\in(A+\gamma A)^\perp$ with that image; a vector of $I_\gamma$ with prescribed coordinates at the
indices $b\ge j$ exists by \cref{prop:links}(c) (the shared vectors at these indices lie in $I_\gamma$ with leading coefficient $1$), and subtracting it corrects the products with the $e^*_b$, which read off the coordinates
at $e_b$. The first columns are orthogonal to both frames and to all columns. \emph{Isometry.} On vectors orthogonal to $A$ and $\gamma A$ the form factors through the map to pairs; by \cref{lem:liouville}(a) with $\gamma'=\gamma$
and (b), the Gram matrix of the lifts plus $N^{-1}$ times the products of their two $z$-forms is the $Q$-Gram matrix of $v$. \emph{The lists.} Order the vectors $e_a$, $\gamma e_a$ ($a<j$), $z$, $\gamma z$ list by list and
add $N$ to the diagonal entries of $z$ and $\gamma z$, which amounts to adjoining $z'$, $z''$ of norm $N$ orthogonal to everything. The two lists are isotropic, and exchanging the two halves of the rows turns the
Gram matrix into a block antidiagonal matrix whose off-diagonal block is the matrix of products of the subframe with its image; its determinant is $\pm\abs{\delta_*}^2$, and $\delta_*\ne0$ by \cref{lem:padded}(c).

(b) Pass to the quotient $(\Lambda\cap I_\gamma^\perp)/(\Lambda\cap I_\gamma)$, which is unimodular, with a basis $B$ and coordinates $q$; the classes of rational vectors orthogonal to $I_\gamma$ have the products of the
vectors (the first columns of a lift being a basis of $I_\gamma$, orthogonal to the others). Choose $e\in\GL(\Ad_f)$ so that the classes of the lifts of $ve$ form a primitive matrix, the lattice of the vectors whose lifts have
integral classes. In the quotient enlarged by $z'$, $z''$ of norm $N$, the classes of the lifts of $ve$ are orthogonal to the two lists, and with the border of (a) their Gram matrix is that of $ve$ for $Q$; the ranks add up,
since the rank of the quotient is twice the number of $a<j$ plus $k+1$. Jacobi's identity, applied in this enlarged lattice to the primitive lattice of the lifted classes of $ve$ and the lattice of the two lists,
compares $\mathrm{lds}(v,e)$ with the norm $(N^{2f}\mu_j)^2$ of the Gram determinant of the lists (by (a) and \cref{lem:padded}(c)); the enlargement contributes the factor $[N]^2=N^{4f}$, and the lattice of the lists has
index $\mathrm{ind}$ in the integral vectors of its span, finite because $\delta_*\ne0$. This gives $\mathrm{lds}(v,e)\cdot\mathrm{ind}^2=\mu_j^2$. The classes of the $e_a$, $\gamma e_a$ with $a<j$ span the classes of all of them. The chain of lattices obtained by adding in turn the classes of the $e_a$, those of the $k_\gamma e_a$, and the integral vectors of their span has indices divisible by
$\mathrm{iri}(k_\gamma)$, $\mathrm{iri}(k_\gamma^{-1})$ and $t(k_\gamma)$ respectively: the first two through the products with the two frames, which are defined on classes, and the third because a vector of the span differs
from the representative of its class by a vector of $I_\gamma$ in the lattice of the $e_a$ and $k_\gamma e_a$. Hence the product $w$ of the three indices divides $\mathrm{ind}$, is nonzero, and
$\mathrm{lds}(v,e)w^2\le\mu_j^2$. \emph{The level.} The lifts $tw$ of $ve$ have integral parts $m$ in $A$; the vectors $x$ for which $m x$ is congruent modulo $A\cap\gamma A$ to $M$ times a vector of $\Lambda\cap A$ form a
submodule of the integral vectors with a basis $B'$ of index dividing the index of $M\,L$ in $L$, where $L$ is the lattice of the isotropic rows of $Z(k_\gamma)$, a copy of $(\Lambda\cap A)/(\Lambda\cap I_\gamma)$ generated up to finite index by
$i$ vectors; so $[\det B']\le M^{2if}$. The lattice of $veB'$ lies in $L^+_\gamma$ (adding vectors of $I_\gamma$ makes the parts in $A$ divisible by $M$), and
\[
  t^2\,\mathrm{lds}(v,eB')=t^2\mathrm{lds}(v,e)[\det B']^2\le\frac{M^{4if}\mu_j^2}{(\mathrm{iri}\,\mathrm{iri}')^2}=\Bigl(\frac{M^{2if}\mathsf g^\sharp}{N^f}\Bigr)^2.
\]
(c) is \cref{lem:link-vectors}(b), (a), (d).
\end{proof}

\begin{lemma}[The cut]\label{lem:cut}
Let $\gamma$ be a term of cell rank $j$ over $(c,s,c')$ with $s$ agreeing, $x\le\mathrm{mn}(s)$ a positive integer, and $G=\max(\mathfrak g_{\rm Nm},N^f\mu_j)$.
\begin{enumerate}[label=\textup{(\alph*)}]
\item A principal minor of order $r\le1+(j-\varrho)_+$ of the Gram matrix of $z^+$, $\Upsilon_\gamma^{-1}z^+$ has $\mathrm{Nm}\le(G/\mu_j)^2$.
\item If the columns of $w$ have integral lifts (they are images of vectors of $\Lambda$ that are adelic combinations of rational vectors orthogonal to $A$ and $\gamma A$), a minor of the matrix of the $z$-forms of $w$
that is a unit has generalised index at least $x$, provided both forms are chosen unless $x=1$.
\item Let the $r$ columns of $y$ lie in the link, $r\le1+(j-\varrho)_+$ and $r=2$ unless $x=1$, let the $z$-forms be injective on their span, and let them take on the basis $yv$ the values they take on a matrix with
integral lifts. Then $x^2\mu_j^2\le\mathrm{lds}(y,v)\,G^2$.
\item Let $m$ have a left inverse, columns in the link on which both $z$-forms vanish, and let $ma$ lie in $L^+_\gamma$. Then the span $V$ of the parts in $X_0$ of $m$ has dimension $\#\text{columns}$, the lattice
$\Lambda_V\cap L_\gamma$ is a lattice of rank $\dim V$, and its squared discriminant is at most $\mathrm{lds}(m,a)$.
\end{enumerate}
\end{lemma}
\lean{New.Filling.HasCellRank.abs_norm_minor_gram_linkVectors_mul_le, New.Filling.HasIntegralLifts.le_genIndex_det_zRowBlockDiagonal, New.Filling.minorNorm_sq_le_linkDiscSq_mul, New.Filling.HasLevelLifts.discriminant_sq_le, New.Filling.zRowBlockDiagonal_mul_linkMatrix_mul_classCoord, New.Filling.isotropicFramePair_mul_classCoord_eq_zero, New.Filling.toRows_classCoord_mul, New.Filling.submoduleIndex_termLattice, New.Filling.toRows₁_mul_eq_fromRows_zero, New.Filling.gram_linkMetric_eq_gram_toRows₁, New.Filling.HasLevelLifts.hasIntegralLifts, New.Filling.HasLevelLifts.mul, New.Filling.HasLevelLifts.submatrix, New.Filling.hasLevelLifts_of_col_mem, New.Filling.exists_basis_mulVec_mem_ideleM_smul, New.Filling.submoduleIndex_ideleM_smul_isotropicRows_dvd, New.Filling.submoduleIndex_ideleM_smul_eq, New.Filling.submoduleIndex_inf, New.Filling.submoduleIndex_inf_of_le, New.Filling.submoduleIndex_inf_comap, New.Filling.isotropicRowsIndex_mul_torsionIndex_ne_zero, New.Filling.isotropicRowsIndex_inv_pos}

\begin{proof}
(a) The minor is $1$, $N$ or the Gram determinant. $1\le(G/\mu_j)^2$ and $\mathrm{Nm}(N)=N^{2f}\le(G/\mu_j)^2$ because $G\ge N^f\mu_j$. For $r=2$, $j=\varrho+1$ and by \cref{lem:link-vectors}(d) the determinant has modulus
$N^2\abs{\mathfrak g}_v/\abs{D_j}_v^2$ at each place, so its norm is $N^{4f}\mathfrak g_{\rm Nm}^2/\mathrm{Nm}(D_j)^2=(\mathfrak g_{\rm Nm}/\mu_j)^2\le(G/\mu_j)^2$. (b) The values of the $z$-forms are the products of the integral lifts $n$ with $z$
and $k_\gamma z$ in the lattice basis, which are integral; if both are chosen, the lifts are orthogonal to $A$ and $k_\gamma A$, so their products with $(f\mid k_\gamma f)$ vanish except in the two rows of $z$ and $k_\gamma z$, and
\cref{lem:meeting-places}(b) bounds $\mathrm{mn}(s)$ by the index of the $2\times2$ determinant. (c) By \cref{lem:link-vectors}(e) the $z$-forms are the products with the pairs $\varepsilon$; choose $r$ of them that pair
nondegenerately with $y$, say $e$. The pairing $e^*Qyv$ is the minor of the $z$-forms, and its index is at least $x$ by (b); Liouville's inequality (\cref{lem:liouville}(c)) gives $x^2\le\mathrm{lds}(y,v)\,\mathrm{Nm}(\det e^*Qe)$, and by (a) and
\cref{lem:link-vectors}(e), $\mathrm{Nm}(\det e^*Qe)\mu_j^2\le G^2$. (d) Where both $z$-forms vanish the first halves of the pairs lie in $X_0$, the two halves have equal Gram matrices, and $Q$ is the standard Gram
matrix of the parts in $X_0$. The integral lifts $n$ of $ma$ show that the columns of the parts $pa$ lie in $L_\gamma$: $n$ is orthogonal to both frames, so the products of its middle part with the frame of $k_\gamma$
are those of its part in $A$, up to sign, and this part is divisible by $M$. Hence $\Lambda_V\cap L_\gamma$ contains the lattice of $pa$ with finite index in $\Lambda_V$, and its discriminant is at most that of $pa$
(\cref{lem:lattice-disc}(b)).
\end{proof}

\begin{lemma}[Splitting a link along a second one]\label{lem:link-splitting}
Let $\gamma$ be as in \cref{lem:cut}, $\gamma'$ any rational element, and let $x=1$ unless $j=\varrho+1$ and the links of $\gamma$ and $\gamma'$ are equal. The dimension of the intersection of the two links exceeds that of the
common link by at most $1+(j-\varrho)_+$, and by exactly $2$ if $j=\varrho+1$ and the links are equal. There are a basis of a complement of the intersection in the link of $\gamma$, of rank the corank $c^+$, and a lattice in
the complement whose basis vectors, moved along the intersection, have integral lifts, with squared discriminant $D$, such that the part of $L_\gamma$ in the common link is a full lattice there, of discriminant $d$, and
\[
  (t(k_\gamma)\,d\,x)^2D\le\mathsf g^{*2}.
\]
\end{lemma}
\lean{New.Filling.nonempty_linkSplitting, New.Filling.LinkSplitting, New.Filling.HasCellRank.finrank_link_inf_le, New.Filling.HasCellRank.finrank_link_inf_eq_add_two, New.Filling.HasLevelLifts.exists_orthogonal_split, New.Filling.IsKernelExtension.card_add_eq_finrank, New.Filling.LinkSplitting.CommonIndex, New.Filling.LinkSplitting.colSpan_common, New.Filling.LinkSplitting.colSpan_compl_le, New.Filling.LinkSplitting.colSpan_fromCols, New.Filling.LinkSplitting.commonCols, New.Filling.LinkSplitting.complCols, New.Filling.LinkSplitting.complLattice, New.Filling.LinkSplitting.complShear, New.Filling.LinkSplitting.discriminant_sq_mul_le, New.Filling.LinkSplitting.eq_zero_of_mulVec_mem, New.Filling.LinkSplitting.fintypeCommonIndex, New.Filling.LinkSplitting.hasIntegralLifts, New.Filling.LinkSplitting.isLatticeOfRank_commonPart, New.Filling.LinkSplitting.ctorIdx, New.Filling.LinkSplitting.mk.inj, New.Filling.LinkSplitting.mk.injEq, New.Filling.LinkSplitting.mk.sizeOf_spec}

\begin{proof}
\emph{Dimensions.} $\dim X_0^\gamma+(j-\varrho)_+=k$, and a vector of the intersection of the links whose $z$-forms vanish is the pair of a vector of the common link; the $z$-forms are two linear forms, of which only one is
independent of the other below the top rank (\cref{lem:link-vectors}(d)). If the links are equal at the top rank, the intersection is the whole link, of dimension $k+1$, and the common link is $X_0^\gamma$, of dimension $k-1$.
\emph{Splitting.} Take a basis of the link of $\gamma$ in three parts: $m$ spans the kernel of the two $z$-forms in the intersection, $m$ and $u$ together span the intersection, and $y$ completes them. Apply
\cref{thm:link-lattice}(b) to this basis: for some $\phi$ the lattice of $(m\mid u\mid y)\phi$ lies in $L^+_\gamma$ with $t^2\,\mathrm{lds}\le(M^{2if}\mathsf g^\sharp/N^f)^2$. Two integral changes of basis make $\phi$ block
triangular, which splits the squared discriminant into the product $\mathrm{lds}(m,a)\,\mathrm{lds}(u,v_0)\,\mathrm{lds}(y,v)$ of three graded pieces, each moved along the earlier pieces by shears and still with integral lifts.
The part $m$: its parts in $X_0$ span the common link, and \cref{lem:cut}(d) gives $d^2\le\mathrm{lds}(m,a)$. The part $u$ has rank $r$ at most $1+(j-\varrho)_+$, with $r=2$ when $x\ne1$, and the $z$-forms are injective on its span;
\cref{lem:cut}(c) gives $x^2\mu_j^2\le\mathrm{lds}(u,v_0)G^2$. The part $y$ has rank $c^+$, by the count of dimensions. Hence
\[
  (t\,d\,x)^2D\,\mu_j^2\le t^2\,\mathrm{lds}(m,a)\,\mathrm{lds}(u,v_0)\,\mathrm{lds}(y,v)\,G^2\le\Bigl(\frac{M^{2if}\mathsf g^\sharp G}{N^f}\Bigr)^2=\mathsf g^{*2}\mu_j^2,
\]
using $\mathsf g^\sharp G/(N^f\mu_j)=\mathsf g^{\max}$ (\cref{lem:padded}(e)).
\end{proof}

\begin{lemma}[The congruence and the index]\label{lem:congruence}
Let $\gamma$, $\gamma'$ be terms over one block $(c,s,c')$.
\begin{enumerate}[label=\textup{(\alph*)}]
\item If $w$, $w'$ have integral lifts for $\gamma$, $\gamma'$, then $w^*(M^{-1}\mathrm{lf})w'$ is integral.
\item If $b'$, $b$ are $i$ columns of $X_0\perp Ez^+$ with the same Gram matrix, then $\mathrm{Nm}(\det(b-b')^*(b-b'))\le4^{2if}\mathrm{Nm}(\det b^*b)$.
\item If $y$, $y'$ have columns in the links of $\gamma$, $\gamma'$, $\mathrm{lf}$ pairs them nondegenerately, and $v^*(M^{-1}y^*\mathrm{lf}\,y')v'$ is integral, then $1\le(2/M)^{4if}\,\mathrm{lds}(y,v)\,\mathrm{lds}(y',v')$.
\item The part of $L_\gamma$ in the common link contains $L_{12}$ with an index dividing $t(k_{\gamma'})$.
\end{enumerate}
\end{lemma}
\lean{New.Filling.HasIntegralLifts.isIntMatrix_linkForm, New.Filling.abs_norm_det_gram_sub_le, New.Filling.one_le_mul_linkDiscSq, New.Filling.LinkSplitting.one_le_mul_linkDiscSq, New.Filling.det_conjTranspose_mul_linkForm_mul_ne_zero, New.Filling.abs_norm_det_smul, New.Filling.submoduleIndex_pairLattice_dvd_torsionIndex, New.Filling.map_commonPart_le, New.Filling.pairLattice_eq_commonPart_inf, New.Filling.frameImage_conjTranspose_mulVec, New.Filling.mulVec_sub_mem_colSpan_isotropicFramePair, New.Filling.exists_sub_one_mul_eq_isotropicFramePair_mul_add, New.Filling.translateByClassReps_mul_classCoord, New.Filling.sub_one_mulVec_mem_colSpan_isotropicFramePair, New.Filling.levelGauge_pos}

\begin{proof}
(a) The element $\delta=\gamma\gamma'^{-1}$ is a term over $(c,1,c)$, so its matrix $k$ in the lattice basis is congruent to $1$ modulo $M$. By \cref{lem:liouville}(a) the pairing of the lifts $n$, $n'$ is
$M^{-1}(n^*Jn'-n^*Jkn')=-n^*J\bigl(M^{-1}(k-1)\bigr)n'$ with the integral Gram matrix $J$ of the lattice, which is integral. (b) The Gram matrices of $b-b'$ and $b+b'$ are positive semidefinite and add up to $4$ times that of
$b$, which is $\Theta+\Theta^*$ for $\Theta$ twice the Gram matrix of $b$; so at each place the first is at most $\Theta+\Theta^*$ minus a positive semidefinite matrix, and its determinant is at most $2^i\abs{\det\Theta}$ (Ostrowski--Taussky,
\cite{OstrowskiTaussky1951}), that is $4^i$ times the determinant of the Gram matrix of $b$; take the product over the places. (c) Let $a$, $b$ be the sources of $y$, $y'$ and $u=\Upsilon_\gamma^{-1}\Upsilon_{\gamma'}$. By
\cref{lem:liouville}(a), (b) the link form pairs $y$ with $y'$ as $a$ with $(1-u)b$ in the positive space, so $M^{-1}a$ and $(1-u)b$ pair integrally after $v$, $v'$ and nondegenerately. Liouville's inequality gives
$1\le[\det v]^2\mathrm{Nm}(\det(M^{-1}a)^*(M^{-1}a))\,\mathrm{Nm}(\det((1-u)b)^*((1-u)b))[\det v']^2$, and with $\mathrm{Nm}(M^{-1})^{2i}=M^{-4if}$ and (b) this is at most $(2/M)^{4if}\mathrm{lds}(y,v)\,\mathrm{lds}(y',v')$.
(d) The common lattice $L_{12}$ is the intersection of $L_{\gamma'}$ with the part of $L_\gamma$ in the common link. A vector $v$ of that part has a lift $y=v-Ma$, $a\in\Lambda\cap A$, orthogonal to the frame of $\gamma$.
Let $m=\gamma'\gamma^{-1}$, which is $\equiv1$ modulo $M$ in the lattice basis. Then $my$ is orthogonal to the frame of $\gamma'$, since $h(my,\gamma'f)=h(\gamma^{-1}y,f)=h(y,\gamma f)=0$; and $(m-1)y$ lies in $A+\gamma'A$, because the
link isometries of $\gamma$ and $\gamma'$ agree on $v$: if $t'\in v+A$ is orthogonal to $A$ and $\gamma'A$, then $\gamma^{-1}y-\gamma'^{-1}t'\in A$ and $(m-1)y=(t'-y)+\gamma'(\gamma^{-1}y-\gamma'^{-1}t')$. Hence
$v=Ma+Mw+my$ with $w=-M^{-1}(m-1)y$ in $\Lambda\cap(A+k_{\gamma'}A)$, and the products of $v$ with the frame of $k_{\gamma'}$ are $M$ times those of $a+w$. The vectors of $k_{\gamma'}(\Lambda\cap A)$ have products $0$ with that
frame ($A$ is totally isotropic and orthogonal to $z$). Let $\pi(x)$ be the vector of the products of $x$ with the frame of $k_{\gamma'}$. The homomorphism $v\mapsto M^{-1}\pi(v)$ modulo $\pi(\Lambda\cap A)$ takes values in
$\pi(\Lambda\cap(A+k_{\gamma'}A))/\pi(\Lambda\cap A)$, a quotient of $\Lambda\cap(A+k_{\gamma'}A)/(\Lambda\cap A+k_{\gamma'}(\Lambda\cap A))$, whose order is $t(k_{\gamma'})$; its kernel is $L_{12}$, by the definition of $L_{\gamma'}$. The
index of $L_{12}$ is the order of the image, a divisor of $t(k_{\gamma'})$. Positivity of the gauges: the indices are finite by \cref{thm:link-lattice}.
\end{proof}

\begin{theorem}[Gain]\label{thm:gain}
Let $\gamma$, $\gamma'$ be terms of cell ranks $j$, $j'$ over one block $(c,s,c')$ with $s$ agreeing, and $x\le\mathrm{mn}(s)$. Then
\[
  \disc(L_{12})^2\le(2/M)^{2c^+f}\,\mathsf g^*(\gamma)\mathsf g^*(\gamma'),
\]
$c^+$ the codimension of the intersection of the two links in each; if $c^+=0$ and $\gamma$ has the top rank, the bound may be divided by $x^2$.
\end{theorem}
\lean{New.Filling.pairDiscriminant_sq_le, New.Filling.LinkSplitting.one_le_mul_linkDiscSq, New.Filling.nonempty_linkSplitting, New.Filling.torsionIndex_mul_discriminant_commonPart_le, New.Filling.submoduleIndex_pairLattice_dvd_torsionIndex, New.Filling.minorNorm_sq_le_linkDiscSq_mul, New.Filling.HasCellRank.abs_norm_minor_gram_linkVectors_mul_le, New.Filling.HasLevelLifts.discriminant_sq_le, New.Filling.HasIntegralLifts.isIntMatrix_linkForm, New.Filling.isLatticeOfRank_and_pairDiscriminant_le, New.Filling.isLatticeOfRank_pairLattice, New.Filling.one_le_pairDiscriminant, New.Filling.pairLattice_comm, New.Filling.pairDiscriminant_comm, New.Filling.IsNear.levelGauge_le, New.Filling.commonLink_comm, New.Filling.linkCorank_comm, New.Filling.linkCorank_eq_zero_iff, New.Filling.le_finrank_commonLink_add, New.Filling.termDisc_eq_zero_iff_of_link_eq, New.Filling.one_le_mul_linkDiscSq, New.Filling.map_commonPart_le, New.Filling.pairLattice_eq_commonPart_inf}

\begin{proof}
Let $X=x$ if $c^+=0$ and $j=\varrho+1$, and $X=1$ otherwise. If $X\ne1$ the links are equal, so $\mathfrak g$ vanishes for both or for neither (the link determines $X_0^\gamma$, which is $X_0$ if $\mathfrak g=0$ and a hyperplane
otherwise, \cref{lem:link-vectors}(d)); hence $\gamma'$ has the top rank as well. \emph{Congruence.} Split both links along their intersection by \cref{lem:link-splitting}, with complements $y$, $y'$ of rank $c^+$ and
squared discriminants $D$, $D'$. A vector of the complement of $\gamma'$ orthogonal for $\mathrm{lf}$ to the complement of $\gamma$ is orthogonal to the whole link of $\gamma$ (it is orthogonal to the intersection as well, since the intersection lies in the
link of $\gamma'$, which is isotropic for $\mathrm{lf}$), and that link is its own orthogonal, so it lies in both links and vanishes; the pairing is nondegenerate, and the terms along the intersection vanish. By \cref{lem:congruence}(a), (c),
$1\le(2/M)^{4c^+f}DD'$. \emph{Cut.} \cref{lem:link-splitting} gives $(t\,d\,X)^2D\le\mathsf g^{*2}$ and $(t'\,d'\,X)^2D'\le\mathsf g^{*\prime2}$, with $d$, $d'$ the discriminants of the parts of $L_\gamma$ and $L_{\gamma'}$ in the common link.
Multiplying and using the congruence,
\[
  (t\,d\,t'\,d'\,X^2)^2\le\bigl((2/M)^{2c^+f}\mathsf g^*\mathsf g^{*\prime}\bigr)^2.
\]
\emph{Index.} By \cref{lem:congruence}(d) and \cref{lem:lattice-disc}(b), $\disc(L_{12})\le d\,t'$ and, exchanging the two terms, $\disc(L_{12})\le d'\,t$. Hence $(\disc(L_{12})X)^2\le t\,d\,t'\,d'\,X^2\le(2/M)^{2c^+f}\mathsf g^*\mathsf g^{*\prime}$.
\end{proof}

\subsection{Counting lattices and frames}\label{ss:frame-count}
A class $c$ has the lattice $\Lambda\cap X_0$ of the vectors $v$ for which $g^{-1}v$ is integral with entry $k_0$ divisible by $\bar\pi$, $g$ the chosen representative of $c$. The \emph{lattice stabiliser cosets} are the cosets, in the
stabiliser of the standard lattice of $X_0$, of the subgroup of the elements whose row $k_0$ is congruent to that of $1$ modulo $\bar\pi$; by admissibility there are at least $(N^k/b)^f$ of them. Write $H=b^{n^3f}$.
\lean{New.Filling.InverseRowsCongruent, New.Filling.LatticeStabilizerCoset, New.Filling.BoundedDiscSublattice, New.Filling.unimodularLattice, New.Filling.RowVanishesModPiN, New.Filling.discWeightIfRowVanishes, New.Filling.subspaceDiscWeight, New.Filling.ClassBoundedDiscLattice, New.Filling.discriminantBound}

\begin{proposition}[Lattices of given corank]\label{prop:lattices-corank}
Let $r+l=k$. The sum of $\disc(L)^{-k}$ over all classes $c$ and all lattices $L\subseteq\Lambda\cap X_0$ of rank $r$ with $\disc L\le H$ is at most $\tau^{l+1}n_XN^{-kf}$.
\end{proposition}
\lean{New.Filling.sum_discriminant_inv_le, New.Filling.sum_classes_le_of_sum_translates_le, New.Filling.inv_mul_mem_of_inverseRowsCongruent, New.Filling.exists_mk_out_eq_embFin_mul, New.Filling.map_le_and_rowVanishesModPiN_of_out_eq, New.Filling.tsum_discriminant_inv_le_of_out_eq, New.Filling.genIndex_piN_eq_pow, New.Filling.mem_smul_intVectors_iff, New.Filling.sub_mem_smul_intVectors_of_values, New.Filling.encard_le_of_vecMul_mem_smul, New.Filling.encard_rowVanishesModPiN_le, New.Filling.exists_saturatedLattice_encard_rowVanishesModPiN_le, New.Filling.mul_inv_pow_le_of_mul_le, New.Filling.one_le_sublatticeDiscriminant_unimodular, New.Filling.sum_discWeightIfRowVanishes_eq_encard_mul, New.Filling.exists_sum_discWeightIfRowVanishes_le, New.Filling.sum_tsum_discWeightIfRowVanishes_le, New.Filling.tailScale_nonneg, New.Filling.add_mul_succ_pow_three_le, New.Filling.pow_mul_discriminantBound_pow_le, New.Filling.pow_mul_mul_rootDiscr_pow_le_tailScale, New.Filling.mul_rootDiscr_pow_le_tailScale_pow_div}

\begin{proof}
\emph{Means over the classes.} Let $A$ be a finite set mapped by $\varphi\mapsto a_\varphi$ into the finite adelic unitary group of $X_0$, and $F\ge0$ a function on the classes. If $\sum_{\varphi\in A}F(\text{class of }ua_\varphi)\le\#A\cdot B$ for every
$u$, then $\sum_cF(c)\le n_XB$: the defect $B-F$, a function on the double cosets, has nonnegative sums over the translates, and its integral over the double coset space, a weighted sum of its values with
the weights of the classes, is translation invariant; as the level fixes it, the integral is $\sum_c(B-F(c))$ up to the common weight. Apply this with the lattice stabiliser cosets, $a_\varphi$ an element of $\varphi$.

\emph{Unimodular lattices.} The chosen representative of the class of $ua_\varphi$ is $\eta uy$ with $\eta$ rational and $y$ in the coset $\varphi$, because the level is the preimage of $K(M)$. If $L$ lies in the lattice of that class, then
$\eta^{-1}L$ lies in $u\hat\OO^\kappa$ and the linear form \emph{row $k_0$ of $y^{-1}u^{-1}$}, the \emph{form} of $\varphi$, vanishes on it modulo $\bar\pi$; the rational isometry $\eta^{-1}$ keeps discriminants
(\cref{lem:lattice-disc}(c)). Hence the sum over the lattices of the class of $ua_\varphi$ is at most the sum over the lattices $L\subseteq u\hat\OO^\kappa$ of rank $r$, $\disc L\le H$, on which the form of $\varphi$ vanishes modulo $\bar\pi$.

\emph{Few forms vanish on a lattice.} Let $L\subseteq u\hat\OO^\kappa$ have rank $r$. By \cref{lem:lattice-disc}(d) it has a finite index $i$ in the saturated lattice $\Lambda_V$ of its span $V$, with basis $mA$ for a basis $m$ of
$\Lambda_V$, $[\det A]=i$, and $\disc L=i\,d(V)$. Extend $u^{-1}m$ to a basis $U$ of $\hat\OO^\kappa$. Integral forms whose values on $mA$ are divisible by an integral idele $d$ and which are pairwise incongruent modulo $d$ are at most
$[\det A]\,[d]^l$ in number: write the values on $mA$ as $dW$; the class of $W$ modulo the rows of $A$, of which there are $[\det A]$, fixes the values on $m$ modulo $d$, and the values on the other $l$ basis vectors have $[d]^l$
classes. Two cosets whose forms are congruent modulo $\bar\pi$ are equal (their inverses have congruent rows $k_0$). With $d=\bar\pi$ and $[\bar\pi]=N^f$, at most $i(N^f)^l$ cosets $\varphi$ have a form vanishing on $L$.

\emph{Summing.} Each such $\varphi$ contributes $(d(V)i)^{-k}$, and for $d(V)$, $i\ge1$ with $d(V)i\le H$ and $r\le k$,
\[
  i\,(d(V)i)^{-k}\le H^{\mu+2}d(V)^{-(k+1)}i^{-(r+\mu)},
\]
since $d(V)\,i^{1+\mu}\le H^{1+\mu}$. Sort the lattices $L$ by their span $V$; with $\mu=2\varrho+4$, \cref{lem:sublattice-zeta}(c) sums over the sublattices of $\Lambda_V$ and \cref{prop:siegel-lemma}(c) over the subspaces
of codimension $l$:
\[
  \sum_\varphi\sum_L\disc(L)^{-k}\le(N^f)^lH^{2\varrho+6}\,(200fk\,\mathrm{rd}_E)^{f(k+1)l}\,2^k.
\]
\emph{Constants.} $b^fH^{2\varrho+6}2^k\le\tau$, as $1+n^3(2\varrho+6)+1\le10^5n^9$; and $N^f(200fk\,\mathrm{rd}_E)^{f(k+1)}\le\tau$, comparing thousandth powers with $(200fk\,\mathrm{rd}_E)^{1000}\le b$ and $N\le b^{60n}$. Hence the right side is
at most $\tau^{l+1}/b^f$. Dividing by the number of cosets, at least $(N^k/b)^f$, the mean is at most $\tau^{l+1}N^{-kf}$, and the first step gives the claim.
\end{proof}

For a term $\gamma$ of cell rank $j$ let $\Xi_0$ be the first $\varrho$ rows of $Z=Z(\gamma)$, $\mathrm{fm}_j(Z)=M\det(P_jP_j^*)$ the \emph{frame modulus}, and let the \emph{shared vectors} be the coordinates in $A$ of the columns $b\ge j$
of $\gamma f\cdot\operatorname{adj}P_j$ (the others set to $0$). A \emph{move} is an element of $U^0_c(M)$. The \emph{part in $X_0$} of the frame of $\gamma$ is the matrix of the coordinates in $X_0$ of the vectors $\gamma f_a$.
\lean{New.Filling.sharedVectors, New.Filling.frameModulus, New.Filling.framePartX₀, New.Filling.IsLift, New.Filling.complementLattice, New.Filling.skewSolution, New.Filling.IsFrameSet}

\begin{lemma}[Short vectors and moves]\label{lem:frame-moves}
Let $\gamma$, $\gamma'$ have cell rank $j$ over blocks with first class $c$.
\begin{enumerate}[label=\textup{(\alph*)}]
\item For $b\ge j$ or $b$ the last index, the combination of $\gamma f$ with the coefficients of column $b$ of $P_j^{-1}$ has no coordinates at the $e^*_a$, and its part in $X_0$ has squared length at most $1/5$ at every
place. A vector of $\Lambda\cap X_0$ of squared length less than $1$ at every place is $0$.
\item If $\gamma$ is near, the entries of $\operatorname{adj}P_j$ are at most $b^{(2+134n)\varrho}$ and those of the shared vectors at most $b^{182n^2}$ at every place; $[MD_j]\le b^{2n^3f}$ and $[\mathrm{fm}_j(Z)]\le b^{2n^3f}$.
\item Moves fix the frame, keep the products, the cell rank, the shared vectors and the blocks. Near terms with the same label whose frames agree after moves are equal.
\item If $Z(\gamma')=Z(\gamma)$ is integral, the shared vectors agree, the frames agree outside the isotropic space and their isotropic rows differ by $\mathrm{fm}_j(Z)X$ with $X$ integral, then a move carries the frame of $\gamma$ to
that of $\gamma'$.
\item For a matrix $\eta$ with columns in $\Lambda\cap X_0$ there is a move subtracting $M\eta\Xi_0$ from the part in $X_0$ of every frame. If $Z(\gamma')=Z(\gamma)$ is integral and the parts in $X_0$ are congruent modulo
$MD_j(\Lambda\cap X_0)$, they differ by $M\eta\Xi_0$ for such an $\eta$, and a move makes them equal.
\item If the lattice $L$ with basis $ya$ lies in $L_\gamma$, it has a rational lift: a matrix $w$ with $wa$ integral such that every vector of $L$ plus $M$ times the corresponding combination of the $e_a$ is orthogonal to $\gamma f$.
Two lifts whose rows are congruent modulo the rows of the Gram matrix of $ya$, a lattice of index $\disc(L)^2$ in the integral rows, are joined by a move. The lattice $\Lambda\cap X_0\cap L^\perp$ has rank $l=k-r$.
\end{enumerate}
\end{lemma}
\lean{New.Filling.HasCellRank.mulVec_eStar_eq_zero_and_sum_indexX₀_le, New.Filling.col_eq_zero_of_norm_lt_one, New.Filling.IsNearTermAt.eq_of_frameImage_eq, New.Filling.norm_adjugate_le, New.Filling.norm_adjugate_le_pow, New.Filling.norm_coeff_le_norm_sum_frameVector, New.Filling.sum_smul_frameVector_eq, New.Filling.norm_basisScale_besselFrame_le, New.Filling.norm_sum_div_basisScale_smul_frameVector_le, New.Filling.IsNear.norm_adjugate_paddedLeading_le, New.Filling.HasCellRank.exists_frameImage_one_mulVec_eq, New.Filling.frameProducts_mul_of_mem_besselU0, New.Filling.HasCellRank.sharedVectors_mul_of_mem_besselU0, New.Filling.HasCellRank.norm_sharedVectors_le, New.Filling.conjTranspose_skewSolution, New.Filling.skewSolution_mul, New.Filling.isIntMatrix_skewSolution, New.Filling.isStarStable_comap_Ô, New.Filling.frameModulus_ne_zero, New.Filling.isIntMatrix_paddedLeading, New.Filling.mul_paddedLeading_eq_submatrix_mul_isotropicRows, New.Filling.mul_adjugate_mul_isotropicRows, New.Filling.mulVec_adjugate_top_eq_zero, New.Filling.inclMatrix_conjTranspose_mul_frameImage, New.Filling.conjTranspose_mul_isotropicRows_eq_neg, New.Filling.frameImage_sub_mul_adjugate_apply_eq_zero, New.Filling.exists_frameImage_mul_eq_add_fromRows, New.Filling.exists_mem_classUnipotent₀_frameImage_mul_eq, New.Filling.exists_framePartX₀_mul_eq_sub_smul, New.Filling.HasCellRank.exists_framePartX₀_sub_eq, New.Filling.HasCellRank.exists_framePartX₀_mul_eq_of_sub_mem, New.Filling.HasCellRank.sum_norm_sq_sub_mul_inv_lt_one, New.Filling.exists_isLift, New.Filling.IsLift.exists_conjTranspose_framePartX₀_mul_eq_neg, New.Filling.IsLift.isIntMatrix, New.Filling.IsLift.orthogonal, New.Filling.exists_conjTranspose_mul_eq_of_row_mem, New.Filling.submoduleIndex_gram_eq_discriminant_sq, New.Filling.submoduleIndex_colLattice_gram_transpose, New.Filling.coe_mem_complementLattice, New.Filling.submoduleIndex_smul_complementLattice, New.Filling.IsNear.genIndex_M_mul_leadingMinor_le, New.Filling.IsNear.genIndex_frameModulus_le, New.Filling.genIndex_algebraMap_le_pow_of_norm_sq_le, New.Filling.frameModulus_mem_integralAdeles_of_isIntegral, New.Filling.frameProducts_castSucc_apply, New.Filling.frameProducts_top_apply, New.Filling.apply_inr_besselFrame_eq, New.Filling.frameMiddle_translateByClassReps, New.Filling.frameMiddle_translateByClassReps_mul_classCoordX₀, New.Filling.col_framePartX₀_sub_mem, New.Filling.exists_mul_eq_of_finAdeleMatrix_mul_eq, New.Filling.exists_isIntMatrix_mul_eq_of_col_mem, New.Filling.colLattice_smul, New.Filling.injective_restrictScalars_mulVecLin, New.Filling.isIntMatrix_mul_finAdeleMatrix_inv_smul, New.Filling.submoduleIndex_colLattice_mul, New.Filling.coe_mem_saturatedLattice, New.Filling.submoduleIndex_smul_saturatedLattice}

\begin{proof}
(a) For $b\ge j$ the vector lies in the span of $f$ (\cref{lem:padded}(b)), where the coordinates at the $e^*_a$ and in $X_0$ vanish. For the last index below $j$ we have $j=\varrho+1$, $P_j=Z$ and $Zx=e_z$; the vector
$y=\gamma fx$ is orthogonal to the $e_a$, so it has no coordinates at the $e^*_a$, and $h(z,y)=1$, so its coordinate along $z$ is $-1/N$. Comparing $h(y,y)=\abs{y_0}^2-1/N$, with $y_0$ the part in $X_0$, with
$h(y,y)=h(fx,fx)=-N\abs{x_z}^2$ gives $N\abs{y_0}^2=1-\abs{Nx_z}^2\le1$ at every place, and $N\ge5$. The squared length of a vector of $\Lambda\cap X_0$ is an algebraic integer; if all its conjugates are less than $1$ it is $0$, and the form is definite on $X_0$.
(b) Entries of an adjugate are minors of order $\varrho$ up to sign, so they are at most $(n^2B)^n$ if the entries are at most $B$; the entries of $Z$ are at most $b^{134n}$ for a near term. For the shared vectors, write the
vector as $\gamma fx=fw$ with $x$ a column of $\operatorname{adj}P_j$. At the boxed point $\gamma^{-1}y$ its length is bounded because $x$ is: with scales of the frame differing by at most $b^{23\varrho}$ and frame vectors of
length at most $b^{37n}$, it is at most $b^{(2+134n)\varrho}b^{1+23\varrho+37n}$ divided by a scale. At the boxed point $y$, a coefficient of a combination of the frame vectors is at most $b^{46n^2}$ times its length:
take $u$ unipotent with $uy$ small; its unitriangular matrix $T$ on the frame and $Tx$ are small, and $x=\operatorname{adj}T\,(Tx)$. The exponents add up to at most $182n^2$. For the indices: an element of $E$ whose squared
modulus is at most $b^e$ at every place has generalised index at most $b^{ef}$; and $\abs{MD_j}^2$, $\abs{\mathrm{fm}_j}^2\le b^{2n^3}$ at every place.
(c) A move $u$ fixes every frame vector, so $Z(u\gamma)=Z(\gamma)$, and the columns $b\ge j$ of $\gamma f\operatorname{adj}P_j$ lie in the span of $f$, which $u$ fixes; $h_c^{-1}uh_c$ is in the principal congruence subgroup, so $u$ is a
term over $(c,1,c)$. If $u\gamma$ and $u'\gamma'$ have equal frames, the reductions of the two terms lie in one coset of $\bar R_0$, and their labels are equal, so they have the same reduction; the classes agree, the terms lie in
one double coset, and both are its chosen element.
(d) With $P=\Xi_0$, $\delta=\det P_j$ and $\alpha$ the first $\varrho$ columns of $\operatorname{adj}P_j$, the difference of the isotropic rows is $\xi=\mathrm{fm}_jX$, and $\xi^*P$ is skew-hermitian because both frames have the Gram matrix of
$f$. The columns $b\ge j$ of $X\operatorname{adj}P_j$ vanish, as the shared vectors agree, and at the top cell rank the last column vanishes too (a matrix $\mu$ with $P^*\mu$ skew-hermitian kills the last column of $\operatorname{adj}Z$);
hence $X\alpha P=\delta X$. The matrix $\sigma=\bar\delta X\alpha-\delta(X\alpha)^*-\alpha^*P^*X\alpha$ is skew-hermitian, integral, and solves $\sigma P=\delta\bar\delta X$. An Eichler element of level $M$ with matrix $\sigma$ and zero part in the lattice of
$c$ fixes $f$ and adds $M\sigma P=\xi$ to the isotropic rows of every frame (\cref{lem:eichler-elements}).
(e) The Eichler map made from $\eta$, orthogonal to the frame, fixes the frame; in the coordinates of $X_0$ it subtracts $M\eta$ times the products with the $e_a$, which are the rows $\Xi_0$. If the difference $v$ of the parts in
$X_0$ is congruent to $0$ modulo $MD_j$, then $vP_j^{-1}=M\eta'$ with columns of $\eta'$ in $\Lambda\cap X_0$; the last column and the columns $b\ge j$ of $vP_j^{-1}$ are differences of two vectors of squared length at most $1/5$ by
(a), so their squared length is less than $1$, and by (a) they vanish. Hence $v=M\eta\Xi_0$ with $\eta$ the remaining columns, and the move of the first statement makes the parts equal.
(f) By the definition of $L_\gamma$ each basis vector has a lift over the finite adeles; the lifts solve a rational linear system, so there is a rational solution near any adelic one. Two lifts give maps $L\to M(\Lambda\cap A)$
whose difference, if each row is an integral combination of the rows of $(ya)^*(ya)$, is given by the products with a row $\eta$ of rational vectors of $L\subseteq\Lambda\cap X_0$ ($\eta=y(y^*y)^{-1}w^*$); by (e) a move
removes it. The rows of the Gram matrix span the image of $L$ in its dual, of index $\disc(L)^2$. The last statement is the rank of an orthogonal complement in a definite space.
\end{proof}

\begin{proposition}[Frames over a lattice]\label{prop:frames-over-lattice}
Let $L\subseteq\Lambda\cap X_0$ be a lattice of corank $l$ with $\disc L\le H$, and $B$ a set of near terms over blocks with first class $c$, with frames distinct modulo moves, each of some cell rank and with $L\subseteq L_\gamma$.
Then $\#B\le b^{2000n^6(l+1)f}$. In particular at most $b^{2000n^6(l+1)f}$ near representatives of a given label over $c$ have $L\subseteq L_\gamma$.
\end{proposition}
\lean{New.Filling.IsFrameSet.encard_le, New.Filling.encard_setOf_isNearTermAt_and_le_termLattice_le, New.Filling.IsFrameSet.encard_le_of_le_termLattice, New.Filling.IsFrameSet.encard_le_of_conjTranspose_mul_eq, New.Filling.IsFrameSet.encard_le_of_framePartX₀_eq, New.Filling.IsFrameSet.encard_le_of_sharedVectors_eq, New.Filling.IsFrameSet.encard_range_frameProducts_le, New.Filling.IsFrameSet.encard_range_sharedVectors_le, New.Filling.IsFrameSet.isIntegral_frameProducts, New.Filling.IsFrameSet.mem_termsOver, New.Filling.IsFrameSet.mono, New.Filling.IsFrameSet.near, New.Filling.IsFrameSet.separated, New.Filling.encard_le_of_colouring, New.Filling.encard_le_of_encard_range_le, New.Filling.isIntegral_sharedVectors}

\begin{proof}
Colour the terms successively; if every set of terms of one colour has at most $x$ elements and there are at most $C$ colours, there are at most $Cx$ terms, and colouring families of vectors of a module $K'$ modulo a module $K$
uses at most $[K':K\cap K']^{\#\text{family}}$ colours. \emph{Cell rank and products.} There are $\varrho+2\le b^{2f}$ cell ranks. The matrices $Z$ have integral entries of modulus at most $b^{134n}$, so by \cref{lem:box-integral}(a) there are
at most $(1+2b^{134n})^{2n^2f}\le b^{270n^3f}$ of them. Fix $j$ and $Z$, with $d=MD_j$ and $d'=\mathrm{fm}_j(Z)$. \emph{First round: the lift of $L$.} Colour the terms by the rows of their lifts of $L$ (\cref{lem:frame-moves}(f)) modulo the
lattice spanned by the Gram matrix of a basis of $L$; there are at most $(\disc L^2)^\varrho$ colours, and in one colour a move gives every frame the lift of one of them. \emph{Second round: the part in $X_0$.} With equal
lifts, the parts in $X_0$ of the frames have the same products with the basis of $L$, so they differ from that of one of them by vectors of $\Lambda\cap X_0\cap L^\perp$, of rank $l$; sort them modulo $d$ times this lattice,
$[d]^{ln}$ classes; in one class the parts are congruent modulo $d(\Lambda\cap X_0)$ and a move makes them equal (\cref{lem:frame-moves}(e)). \emph{Third round: the isotropic part.} The shared vectors are integral with entries at most
$b^{182n^2}$, so they take at most $b^{366n^4f}$ values; with equal shared vectors and equal parts in $X_0$, a term is determined by the isotropic rows of its moved frame modulo $d'$ (\cref{lem:frame-moves}(d)), which have
$[d']^{n\varrho}$ classes. With $\disc L\le H$ and $[d]$, $[d']\le b^{2n^3f}$ (\cref{lem:frame-moves}(b)) the three rounds give at most $b^{(366n^4+2n^3(\varrho+ln+n\varrho))f}$ terms, and
$2+270n^3+366n^4+2n^3(\varrho+ln+n\varrho)\le2000n^6(l+1)$. Near representatives of one label over $c$ have frames distinct modulo moves by \cref{lem:frame-moves}(c).
\end{proof}

\begin{theorem}[Counting]\label{thm:frame-count}
\begin{enumerate}[label=\textup{(\alph*)}]
\item (Sublattices.) Over a ring whose ideals of finite index are principal, $\sum_L[\Lambda:L]^{-(r+\mu)}$ over the sublattices of finite index of a free module of rank $r$ is at most $\zeta(\mu)^r$, a power of the
sum over the ideals (an upper bound in the direction of Hey's formula); for subspaces $T$ of codimension $s$, $\sum_Td(T)^{-(k+1)}\le(200fk\,\mathrm{rd}_E)^{f(k+1)s}$, by Siegel's lemma (a subspace of codimension $s$
is cut out by $s$ lattice vectors whose norms have a bounded product times its discriminant, Bombieri--Vaaler).
\item (Lattices of corank $l$.) $\sum_c\sum_L\operatorname{disc}(L)^{-k}$, over the lattices $L\subseteq\Lambda\cap X_0$ of corank $l$ and discriminant at most $b^{n^3f}$, is at most $\tau^{l+1}n_XN^{-kf}$.
\item (Frames over a lattice.) Near terms with a cell rank, different modulo moves, whose lattices contain a given $L$ of corank $l$ and discriminant at most $b^{n^3f}$, are at most $b^{2000n^6(l+1)f}$
in number: sorted by the matrix of products, the shared vectors, the part in $X_0$ and the part in the isotropic space, each round costs a bounded power.
\item (Pairs.) The sum of $\operatorname{disc}(L_{12})^{-k}$ over near pairs over agreeing blocks with corank $l\le n^3$ and discriminant at most $b^{n^3f}$ is at most $\tau^{l+2}n_p$.
\end{enumerate}
\end{theorem}
\lean{New.Filling.sum_discriminant_inv_le, New.Filling.encard_le_of_vecMul_mem_smul, New.Filling.IsFrameSet.encard_le, New.Filling.IsFrameSet, New.Filling.IsFrameSet.encard_le_of_framePartX₀_eq, New.Filling.IsFrameSet.encard_le_of_conjTranspose_mul_eq, New.Filling.IsFrameSet.encard_le_of_le_termLattice, New.Filling.encard_setOf_isNearTermAt_and_le_termLattice_le, New.Filling.sum_pairDiscriminant_inv_le, New.Filling.siegel_lemma, New.Filling.tsum_sublatticeDiscriminant_inv_le, New.Filling.sublatticeBasis, New.Filling.sublatticeDiscriminant, New.Filling.submoduleZeta, New.Filling.tsum_submoduleIndex_colLattice_inv_pow_le, New.Filling.IsLatticeOfRank.tsum_submoduleIndex_inv_pow_le, New.Filling.discriminant, New.Filling.saturatedLattice, New.Filling.IsLatticeOfRank, New.Filling.jacobi_perp_discriminant, New.Filling.exists_mem_classUnipotent₀_frameImage_mul_eq, New.Filling.IsNearTermAt.eq_of_frameImage_eq, New.Filling.HasCellRank.norm_sharedVectors_le, New.Filling.sharedVectors, New.Filling.frameModulus, New.Filling.exists_framePartX₀_mul_eq_sub_smul, New.Filling.exists_isLift, New.Filling.complementLattice, New.Filling.framePartX₀, New.Filling.HasCellRank.exists_framePartX₀_sub_eq, New.Filling.IsNear.genIndex_frameModulus_le, New.Filling.IsLabelledPair, New.Filling.exists_isLabelledPair_of_isNearPairOverAgreeingBlock, New.Filling.ClassBoundedDiscLattice, New.Filling.boundedDiscNearPairs, New.Filling.sum_pairDiscriminant_inv_le_tsum, New.Filling.tsum_le_tsum_mul_of_encard_le, New.Filling.two_mul_count_exponent_le, New.Filling.BoundedDiscSublattice, New.Filling.unimodularLattice, New.Filling.RowVanishesModPiN, New.Filling.tsum_discriminant_inv_le_of_out_eq, New.Filling.encard_rowVanishesModPiN_le}

\begin{proof}
(a) is \cref{lem:sublattice-zeta}(b), (c) and \cref{prop:siegel-lemma}(c). (b) is \cref{prop:lattices-corank}, and (c) is \cref{prop:frames-over-lattice}.

(d) If $l>k$ there are no lattices of corank $l$ and the sum is empty; let $l\le k$ and $r=k-l$. Fix $c$.

\emph{A pair determines a label and a lattice.} Let $(\gamma,\gamma')$ be a near pair over $(c,s,c')$ with $s$ agreeing. Both terms lie over the same block, so they have the same label $t$, the component of $s$ in its coset of $\bar R_0$, and both are near terms at $t$.
Their common lattice $L_{12}$ is cut out of $L_\gamma\cap L_{\gamma'}$, so it lies in the lattices of both terms and in $\Lambda\cap X_0$, and it is a lattice of rank the dimension of the common link, here $r$; by hypothesis $\disc L_{12}\le b^{n^3f}$.
Hence the map that sends $(c',\gamma,\gamma')$ to the label $t$ and the lattice $L=L_{12}$ lands in the set of pairs $(t,L)$ with $L\subseteq\Lambda\cap X_0$ of rank $r$ and $\disc L\le b^{n^3f}$, and the weight $\disc(L_{12})^{-k}$ of the pair is the weight $\disc(L)^{-k}$ of its image.

\emph{The fibres.} For a pair $(t,L)$ let $T(t,L)$ be the set of near representatives $(c',\gamma)$ of the label $t$ over $c$ with $L\subseteq L_\gamma$. If $(c',\gamma,\gamma')$ lies over $(t,L)$, both $(c',\gamma)$ and $(c',\gamma')$ lie in $T(t,L)$, and the triple is determined by these two
elements. Hence the fibre over $(t,L)$ has at most $\#T(t,L)^2$ elements, and by (c), as $L$ has corank $l$ and discriminant at most $b^{n^3f}$,
\[
  \#T(t,L)^2\le\bigl(b^{2000n^6(l+1)f}\bigr)^2\le b^{10^5n^9f}\le\tau,
\]
since $l+1\le2n^3$ gives $2\cdot2000n^6(l+1)\le8000n^9$, and every power $b^e$ with $e\le10^5n^9f$ is at most $\tau$.

\emph{The sum.} If every element of a set is sent to an element of the same weight, and each fibre has at most $x$ elements, the sum of the weights is at most $x$ times the sum over the images. Therefore
\[
  \sum_{c'}\sum_{\text{pairs over }(c,c')}\disc(L_{12})^{-k}\le\tau\cdot\#\bar R_0\cdot\sum_L\disc(L)^{-k},
\]
the last sum over the lattices $L\subseteq\Lambda\cap X_0$ of rank $r$ with $\disc L\le b^{n^3f}$. Summing over $c$ and applying (b),
\[
  \sum_c\sum_{c'}\sum_{\text{pairs}}\disc(L_{12})^{-k}\le\#\bar R_0\,\tau\,\tau^{l+1}n_XN^{-kf}=\tau^{l+2}n_p,
\]
with $n_p=\#\bar R_0\,n_XN^{-kf}$.
\end{proof}

\subsection{The remaining terms}\label{ss:remaining-sum}

Two terms lie in the same \emph{link class} if $\mathfrak g$ vanishes for neither and their links are equal.
\lean{New.Filling.InSameLinkClass, New.Filling.Piece, New.Filling.piece, New.Filling.piece_eq_some_iff, New.Filling.IsStable.inter_piece, New.Filling.IsNearPairOverAgreeingBlock.mem_boundedDiscNearPairs, New.Filling.IsNearPairOverAgreeingBlock.ofReal_termSize_mul_le, New.Filling.IsNearPairOverAgreeingBlock.one_le_pairDiscriminant, New.Filling.boundedDiscNearPairs, New.Filling.eq_of_mem_termsOver, New.Filling.isNearPairOverAgreeingBlock_of_mem_nearRepresentativesOver, New.Filling.mem_of_mem_nearRepresentativesOver, New.Filling.mem_termsOver_of_mem_nearRepresentativesOver, New.Filling.IsLabelledPair, New.Filling.IsLabelledPair.isLatticeOfRank, New.Filling.IsLabelledPair.isNearTermAt_left, New.Filling.IsLabelledPair.isNearTermAt_right, New.Filling.IsLabelledPair.le_classLatticeX₀, New.Filling.IsLabelledPair.le_termLattice_left, New.Filling.IsLabelledPair.le_termLattice_right}

\begin{lemma}[Pairs of near terms]\label{lem:pair-sum}
\begin{enumerate}[label=\textup{(\alph*)}]
\item For a near pair over $(c,c')$, the corank $l$ of the common lattice is at most $n^3$ and $\disc L_{12}\le H$.
\item If $(\disc L_{12}\,\eta)^2\le\mathsf g^*(\gamma)\mathsf g^*(\gamma')$ with $\eta>0$, both cell second moments are at most $D$, and $\tau^{l+6}D\eta^{-k}\le B$, then
$\mathrm{size}(\gamma)\mathrm{size}(\gamma')\le B\disc(L_{12})^{-k}/(\tau^{l+2}n_p)$.
\item If every pair of a family of near pairs, one family for each two classes, satisfies the hypotheses of (b) with some $\eta$ and $D$ and a common $B\ge0$, then the products of the sizes over all the families add up to
at most $\tau B$.
\item For a stable $T\subseteq\Gamma^*$, the squared norm of the cut block of $T$ over $(c,s,c')$ is at most the sum over the links $\Lambda$ of the squared norms of the blocks of the terms of $T$ with $\mathfrak g\ne0$ and link $\Lambda$, plus
$\nu^2$ times the sum of $\mathrm{size}(\gamma)\mathrm{size}(\gamma')$ over the pairs of near representatives of $T$ over the block that do not lie in the same link class.
\end{enumerate}
\end{lemma}
\lean{New.Filling.IsNearPairOverAgreeingBlock.corank_le_and_pairDiscriminant_le, New.Filling.cube_add_one_le_tailScale, New.Filling.exists_isLabelledPair_of_isNearPairOverAgreeingBlock, New.Filling.sum_pairDiscriminant_inv_le_tsum, New.Filling.tsum_le_tsum_mul_of_encard_le, New.Filling.two_mul_count_exponent_le, New.Filling.enorm_amplitudeSum_sq_le, New.Filling.enorm_amplitudeSum_mul_le}

\begin{proof}
(a) Let $i=c^+$ and $d=\disc L_{12}\ge1$. By \cref{thm:gain} and \cref{lem:orbit-bound}(a), $d^2\le(2/M)^{2if}H^2$, and $2/M\le b^{-18}$ because $2b\le b^{19}\le M$; so $b^{18if}d\le H$, which gives $d\le H$ and $18i\le n^3$. The corank of the common
lattice is at most $i+(j-\varrho)_+\le i+1$ (\cref{lem:link-splitting}), and $i+1\le n^3$.
(b) By definition $\mathrm{size}^2\le\tau^4D/(n_p\mathsf g^{*k})$, so $\mathrm{size}(\gamma)\mathrm{size}(\gamma')\le\tau^4D/(n_p(\mathsf g^*\mathsf g^{*\prime})^{k/2})\le\tau^4D\eta^{-k}d^{-k}/n_p$, and $\tau^4D\eta^{-k}\le B\tau^{-(l+2)}$.
(c) Sort the pairs by the corank $l\le n^3$ of their common lattice. By (b) and \cref{thm:frame-count}(d) the pairs of corank $l$ contribute at most $B\tau^{-(l+2)}n_p^{-1}\cdot\tau^{l+2}n_p=B$, and there are at most $n^3+1\le\tau$ coranks.
(d) Split $T$ into its pieces: one for each link $\Lambda$, consisting of the terms with $\mathfrak g\ne0$ and link $\Lambda$, and one for the terms with $\mathfrak g=0$. The pieces are stable, since $\mathfrak g$ and the link are invariant
(\cref{prop:links}(a)). By the triangle inequality the squared norm of the block is at most the sum over pairs of pieces of the products of their norms. The squares of the pieces with $\mathfrak g\ne0$ remain. Every other
product is at most $\nu^2$ times the sum of the products of the sizes of the near representatives of the two pieces (\cref{thm:block-size}(c) for each piece), and a pair of near representatives determines its two pieces;
such a pair does not lie in the same link class, since the pieces differ or one of them is the piece of $\mathfrak g=0$.
\end{proof}

\begin{theorem}[Remaining terms]\label{thm:remaining}
Let the data be admissible and $E$ contain a primitive $5040$th root of unity.
\begin{enumerate}[label=\textup{(\alph*)}]
\item Pairs of near terms over agreeing blocks whose links differ, or for one of which $\mathfrak g=0$, contribute $\sum\mathrm{size}(\gamma)\mathrm{size}(\gamma')\le b^{-kf}\mathcal D_{\rm id}$.
\item Pairs with equal links and $\mathfrak g\ne0$ over blocks with $\mathrm{mn}(s)\ge z$ contribute at most $(b^{2f}/z)^k\mathcal D_{\rm main}$.
\item If the terms with $\mathfrak g\ne0$, taken link by link, have link second moment at most $w\nu^2$ over $\Omega$, the second moment over $\Omega$ of $\Gamma^*$ for the cut function is at most
$w+b^{-kf}\mathcal D_{\rm id}$; and if $\mathrm{mn}(s)\ge z$ on $\Omega$, the link second moment over $\Omega$ is at most $(b^{2f}/z)^k\mathcal D_{\rm main}\nu^2$.
\end{enumerate}
\end{theorem}
\lean{New.Filling.sum_termSize_mul_not_inSameLinkClass_le, New.Filling.sum_termSize_mul_inSameLinkClass_le, New.Filling.secondMoment_remainingTerms_le, New.Filling.linkSecondMoment_le_of_le_meetNorm, New.Filling.sum_termSize_mul_le, New.Filling.InSameLinkClass, New.Filling.termSize_sq_le, New.Filling.IsNearPairOverAgreeingBlock, New.Filling.IsNearPairOverAgreeingBlock.termSize_mul_le, New.Filling.IsNearPairOverAgreeingBlock.pairDiscriminant_sq_le, New.Filling.IsNearPairOverAgreeingBlock.corank_le_and_pairDiscriminant_le, New.Filling.cube_add_one_le_tailScale, New.Filling.cellRankSecondMoment_le_mainSecondMoment, New.Filling.pow_mul_mainSecondMoment_div_le, New.Filling.Piece, New.Filling.piece, New.Filling.enorm_amplitudeSum_mul_le, New.Filling.enorm_amplitudeSum_sq_le}

\begin{proof}
(a) Apply \cref{lem:pair-sum}(c) with $B=\tau^{-10}\mathcal D_{\rm id}$; then $\tau B=\tau^{-9}\mathcal D_{\rm id}\le\tau^{-4}\mathcal D_{\rm id}=b^{-kf}\mathcal D_{\rm id}$. Let $(\gamma,\gamma')$ be a pair as in (a), of cell ranks $j$, $j'$, with $i=c^+$. By \cref{lem:link-splitting}
the corank of the common lattice is at most $i+(j-\varrho)_+$. \emph{Equal links.} If $i=0$ the links are equal, so $\mathfrak g$ vanishes for both or for neither (\cref{lem:link-vectors}(d)); as the pair is not in one link class,
it vanishes for both, both terms lie below the top cell rank, the corank is $0$, and $\mathcal D_j=\mathcal D_{j'}=\tau^{-16}\mathcal D_{\rm id}$. Take $\eta=1$, $D=\tau^{-16}\mathcal D_{\rm id}$: \cref{thm:gain} gives
$\disc(L_{12})^2\le\mathsf g^*\mathsf g^{*\prime}$, and $\tau^6D=\tau^{-10}\mathcal D_{\rm id}=B$. \emph{Different links.} If $i\ge1$, take $\eta=(M/2)^{if}$ and $D=\mathcal D_{\rm main}\ge\mathcal D_j,\mathcal D_{j'}$ (\cref{lem:orbit-bound}(b)); by \cref{thm:gain}
$(\disc(L_{12})\eta)^2\le\mathsf g^*\mathsf g^{*\prime}$, and the corank $l$ is at most $i+1$. Put $m=kf$ and $x=\tau(2/M)^m$; since $2b\le M$, $x\le\tau^4(2/M)^m=(2b/M)^m\le1$. Then, with $\mathcal D_{\rm main}\le(M/b^6)^m\mathcal D_{\rm id}$,
\begin{align*}
  \tau^{l+6}\mathcal D_{\rm main}\eta^{-k}&\le\tau^7x^i\mathcal D_{\rm main}\le\tau^7x\,(M/b^6)^m\mathcal D_{\rm id}=\tau^8(2/b^6)^m\mathcal D_{\rm id}\\
  &\le\tau^8(b/b^6)^m\mathcal D_{\rm id}=\tau^{-12}\mathcal D_{\rm id}\le B.
\end{align*}
(b) Here the links are equal and $\mathfrak g\ne0$, so both terms have the top cell rank, $c^+=0$, and \cref{thm:gain} with $x=z\le\mathrm{mn}(s)$ gives $(\disc(L_{12})z)^2\le\mathsf g^*\mathsf g^{*\prime}$. The corank is at most $1$, and the
cell second moments are $\mathcal D_{\rm main}$. Take $\eta=z$, $D=\mathcal D_{\rm main}$ and $B=\tau^7\mathcal D_{\rm main}z^{-k}$ in \cref{lem:pair-sum}(c): the total is at most $\tau^8\mathcal D_{\rm main}z^{-k}=(b^{2f}/z)^k\mathcal D_{\rm main}$, as
$\tau^8=b^{2kf}$.
(c) By \cref{lem:pair-sum}(d), summing over $c$, $c'$ and $s\in\Omega$, the second moment of $\Gamma^*$ is at most the link second moment of the terms with $\mathfrak g\ne0$, divided by $\nu^2$, plus the sum of the products of
the sizes over the pairs of near representatives that do not lie in one link class; a pair determines its $s$, so each such pair occurs once, and by (a) their sum is at most $b^{-kf}\mathcal D_{\rm id}$. For the second
statement apply \cref{lem:pair-sum}(d) to the stable sets of the terms with $\mathfrak g\ne0$ and a fixed link: the square of the block of one link is at most $\nu^2$ times the sum over pairs of its near representatives, which lie in the
same link class; a pair determines its $s$ and its link, and (b) bounds the sum.
\end{proof}
